\global\otCorrigesfalse\global\otSouplefalse\gdef\otAccent{otNavy}\pcAnnexe{Dix brevets blancs}{Volets physique-chimie · Série générale · Préparation 2027}\phantomsection\label{brevets}\pdfbookmark[0]{Dix brevets blancs}{bm-brevets}
\begin{pcEncadre}{method}{Un entraînement au volet de physique-chimie}Chaque sujet est prévu pour 30 minutes et noté sur 10 : 9 points scientifiques et 1 point pour la langue. L’épreuve officielle de sciences dure une heure et comporte deux disciplines retenues parmi trois. Les sujets de ce livre préparent uniquement le volet physique-chimie. Ce sont des créations OptiTAB, sans statut de sujet officiel.\end{pcEncadre}
\pcSousTitre{Conditions de travail}
Prévoir une copie, un brouillon et une calculatrice de collège pour ces entraînements. À l’examen, suivre les consignes du sujet reçu. Chercher sans le cours ; réserver trois minutes à la relecture. Les démarches pertinentes, même incomplètes, peuvent rapporter les points de méthode explicités dans les corrigés.
\pcSousTitre{Le point de langue}
Les quatre dimensions sont appréciées sur la copie entière : orthographe, construction des phrases, vocabulaire scientifique et organisation de la réflexion. Pour l’autoévaluation proposée ici, chacune vaut au plus 0,25 point. Le barème scientifique est séparé. Il ne s’agit pas d’une grille ministérielle imposant cette subdivision.
\pcSousTitre{Parcours conseillé}
Sujets 1 à 3 : commencer après les chapitres correspondants. Sujets 4 à 7 : croiser les thèmes. Sujets 8 à 10 : composer en temps limité sans aide. Tous les documents et nombres d’étude sont originaux ou explicitement annoncés comme références du cours.
\pcRepere{bb-1}{Sujet 01 · Se déplacer en limitant les rejets}
\pcRepere{bb-2}{Sujet 02 · Contrôler une eau au laboratoire}
\pcRepere{bb-3}{Sujet 03 · Éclairer un refuge}
\pcRepere{bb-4}{Sujet 04 · Préparer une mission lunaire}
\pcRepere{bb-5}{Sujet 05 · Entretenir une structure métallique}
\pcRepere{bb-6}{Sujet 06 · Mesurer sous l’eau et préserver l’audition}
\pcRepere{bb-7}{Sujet 07 · Recycler et identifier des matériaux}
\pcRepere{bb-8}{Sujet 08 · Produire et économiser l’électricité}
\pcRepere{bb-9}{Sujet 09 · Observer un système planétaire}
\pcRepere{bb-10}{Sujet 10 · Mener une enquête scientifique}
\pcPage{Brevet blanc 01}{Se déplacer en limitant les rejets}
\phantomsection\label{bb-1}
\pdfbookmark[1]{01 · Se déplacer en limitant les rejets}{bm-bb-1}
\noindent\textbf{Volet PC · 30 min · /10}\hfill Partie 1/2 · 4 points scientifiques
\par\vspace{2mm}Les essais et démarches sont pris en compte. Répondre sur une copie séparée.\par
\begin{pcEncadre}{definition}{Un vélo en mouvement}Un ensemble vélo-cycliste de masse $80\unite{kg}$ parcourt une portion rectiligne de 200 m en 40 s. Le graphique fournit les positions relevées. On donne $E_c=\tfrac12 mv^2$, avec les unités SI. Les valeurs sont des données d’étude.\end{pcEncadre}
\pcFigureEx{figures/bb01-velo.pdf}{}
\par\addvspace{3mm}\noindent\textbf{Question 1 · 1 point}\quad Décrire le mouvement sur cette portion dans le référentiel de la route.\par
\par\addvspace{3mm}\noindent\textbf{Question 2 · 1 point}\quad Calculer la vitesse moyenne en m/s.\par
\par\addvspace{3mm}\noindent\textbf{Question 3 · 2 points}\quad Calculer l’énergie cinétique correspondant à cette vitesse et prévoir sa valeur si la vitesse double.\par
\vfill\pcRepere{cor-bb-1}{Corrigé, barème et remédiation}
\pcPage{Brevet blanc 01}{Se déplacer en limitant les rejets}
\noindent\textbf{Volet PC · 30 min · /10}\hfill Partie 2/2 · 5 points scientifiques
\par\vspace{2mm}Les essais et démarches sont pris en compte. Répondre sur une copie séparée.\par
\begin{pcEncadre}{definition}{Deux solutions énergétiques}Pour une comparaison fictive sur le même trajet : solution A, 20 kWh d’électricité avec une émission de cycle de vie fournie de 50 g de CO2 par kWh ; solution B, combustion émettant 4 000 g de CO2. La réaction modèle du combustible est $\ce{CH4 + 2 O2 -> CO2 + 2 H2O}$. Le tableau ne décrit que l’indicateur CO2 fourni, pas tous les impacts.\end{pcEncadre}
\pcFigureEx{figures/rev-bb1-comparaison.pdf}{Données fictives pour le même trajet et le seul indicateur fourni.}
\par\addvspace{3mm}\noindent\textbf{Question 4 · 1 point}\quad Nommer le comburant de la réaction fournie.\par
\par\addvspace{3mm}\noindent\textbf{Question 5 · 1 point}\quad Compter les atomes d’oxygène de chaque côté de l’équation.\par
\par\addvspace{3mm}\noindent\textbf{Question 6 · 1 point}\quad Calculer la masse de CO2 associée à A en kg.\par
\par\addvspace{3mm}\noindent\textbf{Question 7 · 2 points}\quad En quatre à six lignes, comparer A et B sur l’indicateur donné et préciser une limite environnementale de la conclusion.\par
\vfill\pcRepere{cor-bb-1}{Corrigé, barème et remédiation}
{\small 1 point de langue sur l’ensemble de la copie : voir la grille d’autoévaluation.}
\pcPage{Brevet blanc 02}{Contrôler une eau au laboratoire}
\phantomsection\label{bb-2}
\pdfbookmark[1]{02 · Contrôler une eau au laboratoire}{bm-bb-2}
\noindent\textbf{Volet PC · 30 min · /10}\hfill Partie 1/2 · 4 points scientifiques
\par\vspace{2mm}Les essais et démarches sont pris en compte. Répondre sur une copie séparée.\par
\begin{pcEncadre}{definition}{Masse volumique et mélange}Un récipient vide a une masse de 32,0 g. Avec 50,0 mL d’un liquide, sa masse est 82,0 g. Le liquide est transparent et contient des ions dissous. À température de l’étude, la masse volumique de l’eau pure est voisine de 1,00 g/mL.\end{pcEncadre}
\pcFigureEx{figures/rev-bb2-masse-ex.pdf}{Distinguer récipient rempli et liquide seul.}
\par\addvspace{3mm}\noindent\textbf{Question 1 · 1 point}\quad Calculer la masse de liquide.\par
\par\addvspace{3mm}\noindent\textbf{Question 2 · 1 point}\quad Déterminer sa masse volumique en g/mL.\par
\par\addvspace{3mm}\noindent\textbf{Question 3 · 2 points}\quad Peut-on affirmer que ce liquide est de l’eau pure et potable ? Donner deux arguments distincts.\par
\vfill\pcRepere{cor-bb-2}{Corrigé, barème et remédiation}
\pcPage{Brevet blanc 02}{Contrôler une eau au laboratoire}
\noindent\textbf{Volet PC · 30 min · /10}\hfill Partie 2/2 · 5 points scientifiques
\par\vspace{2mm}Les essais et démarches sont pris en compte. Répondre sur une copie séparée.\par
\begin{pcEncadre}{definition}{Les résultats d’analyse}Le pH vaut 6. Une banque de tests indique : nitrate d’argent + ions chlorure → précipité blanc ; soude + ions cuivre(II) → précipité bleu. Sur deux prélèvements, on obtient un précipité blanc au premier test et aucun bleu au second. Pour Cl, Z = 17.\end{pcEncadre}
\pcFigureEx{figures/rev-bb2-tests-ex.pdf}{Utiliser les observations et la banque de tests du document.}
\par\addvspace{3mm}\noindent\textbf{Question 4 · 1 point}\quad Classer le liquide selon son pH.\par
\par\addvspace{3mm}\noindent\textbf{Question 5 · 1 point}\quad Quel ion est mis en évidence par le précipité blanc ?\par
\par\addvspace{3mm}\noindent\textbf{Question 6 · 1 point}\quad Combien d’électrons possède cet ion monoatomique ?\par
\par\addvspace{3mm}\noindent\textbf{Question 7 · 2 points}\quad En quatre à six lignes, proposer les précautions nécessaires pour comparer les tests et préciser la portée du résultat négatif au cuivre.\par
\vfill\pcRepere{cor-bb-2}{Corrigé, barème et remédiation}
{\small 1 point de langue sur l’ensemble de la copie : voir la grille d’autoévaluation.}
\pcPage{Brevet blanc 03}{Éclairer un refuge}
\phantomsection\label{bb-3}
\pdfbookmark[1]{03 · Éclairer un refuge}{bm-bb-3}
\noindent\textbf{Volet PC · 30 min · /10}\hfill Partie 1/2 · 4 points scientifiques
\par\vspace{2mm}Les essais et démarches sont pris en compte. Répondre sur une copie séparée.\par
\begin{pcEncadre}{definition}{Le circuit d’éclairage}Une source continue idéale de 12 V alimente deux lampes en dérivation. La première est traversée par 0,30 A, la seconde par 0,20 A. Chaque lampe possède son interrupteur en série dans sa branche.\end{pcEncadre}
\pcFigureEx{figures/rev-bb3-ex.pdf}{Deux commandes indépendantes ; interrupteurs fermés pour les mesures.}
\par\addvspace{3mm}\noindent\textbf{Question 1 · 1 point}\quad Donner la tension aux bornes de chaque lampe.\par
\par\addvspace{3mm}\noindent\textbf{Question 2 · 1 point}\quad Calculer l’intensité fournie par la source.\par
\par\addvspace{3mm}\noindent\textbf{Question 3 · 2 points}\quad Décrire où brancher un voltmètre pour la première lampe et un ampèremètre pour la seconde.\par
\vfill\pcRepere{cor-bb-3}{Corrigé, barème et remédiation}
\pcPage{Brevet blanc 03}{Éclairer un refuge}
\noindent\textbf{Volet PC · 30 min · /10}\hfill Partie 2/2 · 5 points scientifiques
\par\vspace{2mm}Les essais et démarches sont pris en compte. Répondre sur une copie séparée.\par
\begin{pcEncadre}{definition}{Une réserve limitée}La puissance totale des lampes est constante et vaut 6,0 W. La batterie peut fournir au plus 36 Wh d’énergie utilisable. On néglige les autres consommations. On donne $E=P\Delta t$ et $1\unite{Wh}=3\,600\unite J$.\end{pcEncadre}
\par\addvspace{3mm}\noindent\textbf{Question 4 · 1 point}\quad Calculer l’énergie utilisée par les deux lampes en 4 h.\par
\par\addvspace{3mm}\noindent\textbf{Question 5 · 1 point}\quad Convertir cette énergie en joules.\par
\par\addvspace{3mm}\noindent\textbf{Question 6 · 1 point}\quad Déterminer la durée maximale théorique à 6,0 W.\par
\par\addvspace{3mm}\noindent\textbf{Question 7 · 2 points}\quad En quatre à six lignes, expliquer pourquoi une durée réelle peut être plus courte et proposer une économie d’énergie justifiée.\par
\vfill\pcRepere{cor-bb-3}{Corrigé, barème et remédiation}
{\small 1 point de langue sur l’ensemble de la copie : voir la grille d’autoévaluation.}
\pcPage{Brevet blanc 04}{Préparer une mission lunaire}
\phantomsection\label{bb-4}
\pdfbookmark[1]{04 · Préparer une mission lunaire}{bm-bb-4}
\noindent\textbf{Volet PC · 30 min · /10}\hfill Partie 1/2 · 4 points scientifiques
\par\vspace{2mm}Les essais et démarches sont pris en compte. Répondre sur une copie séparée.\par
\begin{pcEncadre}{definition}{Peser un équipement}Un équipement de masse 12 kg est transporté de la Terre à la Lune. On donne $g_T=10\unite{N/kg}$, $g_L=1{,}6\unite{N/kg}$ et $P=mg$. L’équipement est posé immobile sur un sol horizontal.\end{pcEncadre}
\pcFigureEx{figures/b09-poids.pdf}{}
\par\addvspace{3mm}\noindent\textbf{Question 1 · 1 point}\quad Donner sa masse sur la Lune.\par
\par\addvspace{3mm}\noindent\textbf{Question 2 · 1 point}\quad Calculer son poids sur la Lune.\par
\par\addvspace{3mm}\noindent\textbf{Question 3 · 2 points}\quad Décrire les deux forces du modèle d’équilibre sur le sol lunaire, avec leur direction, sens et valeur.\par
\vfill\pcRepere{cor-bb-4}{Corrigé, barème et remédiation}
\pcPage{Brevet blanc 04}{Préparer une mission lunaire}
\noindent\textbf{Volet PC · 30 min · /10}\hfill Partie 2/2 · 5 points scientifiques
\par\vspace{2mm}Les essais et démarches sont pris en compte. Répondre sur une copie séparée.\par
\begin{pcEncadre}{definition}{Communiquer avec la Terre}La distance Terre-Lune est prise égale à $3{,}84\times10^8\unite m$. Un signal radio voyage à $3{,}00\times10^8\unite{m/s}$ dans le vide. On néglige tous les retards électroniques. La relation de gravitation fournie est $F=Gm_A m_B/d^2$.\end{pcEncadre}
\pcFigureEx{figures/bb04-radio.pdf}{}
\par\addvspace{3mm}\noindent\textbf{Question 4 · 1 point}\quad Le signal radio est-il sonore ou électromagnétique ?\par
\par\addvspace{3mm}\noindent\textbf{Question 5 · 1 point}\quad Calculer la durée d’un trajet simple.\par
\par\addvspace{3mm}\noindent\textbf{Question 6 · 1 point}\quad Calculer la durée minimale d’un aller-retour.\par
\par\addvspace{3mm}\noindent\textbf{Question 7 · 2 points}\quad En quatre à six lignes, expliquer pourquoi une station en orbite reste soumise à la gravitation et prévoir l’effet d’une distance centrale doublée sur F, à masses fixées.\par
\vfill\pcRepere{cor-bb-4}{Corrigé, barème et remédiation}
{\small 1 point de langue sur l’ensemble de la copie : voir la grille d’autoévaluation.}
\pcPage{Brevet blanc 05}{Entretenir une structure métallique}
\phantomsection\label{bb-5}
\pdfbookmark[1]{05 · Entretenir une structure métallique}{bm-bb-5}
\noindent\textbf{Volet PC · 30 min · /10}\hfill Partie 1/2 · 4 points scientifiques
\par\vspace{2mm}Les essais et démarches sont pris en compte. Répondre sur une copie séparée.\par
\begin{pcEncadre}{definition}{Comparer des protections}Quatre clous identiques sont suivis pendant sept jours. A : air sec, pas de rouille ; B : eau au contact de l’air, rouille ; C : eau salée au contact de l’air, rouille ; D : eau désaérée maintenue sans dioxygène, pas de rouille. Les autres paramètres sont comparables.\end{pcEncadre}
\pcFigureEx{figures/b05-corrosion.pdf}{}
\par\addvspace{3mm}\noindent\textbf{Question 1 · 1 point}\quad La corrosion est-elle une transformation physique ou chimique ?\par
\par\addvspace{3mm}\noindent\textbf{Question 2 · 1 point}\quad Quel facteur est testé en comparant B et D ?\par
\par\addvspace{3mm}\noindent\textbf{Question 3 · 2 points}\quad Ces seuls résultats permettent-ils d’affirmer que le sel double la vitesse de corrosion ? Justifier et proposer une mesure supplémentaire.\par
\vfill\pcRepere{cor-bb-5}{Corrigé, barème et remédiation}
\pcPage{Brevet blanc 05}{Entretenir une structure métallique}
\noindent\textbf{Volet PC · 30 min · /10}\hfill Partie 2/2 · 5 points scientifiques
\par\vspace{2mm}Les essais et démarches sont pris en compte. Répondre sur une copie séparée.\par
\begin{pcEncadre}{definition}{Le fer dans une solution acide}La réaction fournie est $\ce{Fe + 2 H+ -> Fe^{2+} + H2}$. Le fer a Z = 26. Après réaction, un test encadré identifie le gaz ; un précipité vert avec la soude identifie Fe2+. Le pH passe de 2 à 4.\end{pcEncadre}
\par\addvspace{3mm}\noindent\textbf{Question 4 · 1 point}\quad Nommer le gaz produit.\par
\par\addvspace{3mm}\noindent\textbf{Question 5 · 1 point}\quad Déterminer le nombre d’électrons dans Fe2+.\par
\par\addvspace{3mm}\noindent\textbf{Question 6 · 1 point}\quad Le liquide final est-il neutre ? Justifier.\par
\par\addvspace{3mm}\noindent\textbf{Question 7 · 2 points}\quad En quatre à six lignes, relier les observations à la réaction fournie, puis distinguer cette réaction de la corrosion habituelle du fer dans l’air humide.\par
\vfill\pcRepere{cor-bb-5}{Corrigé, barème et remédiation}
{\small 1 point de langue sur l’ensemble de la copie : voir la grille d’autoévaluation.}
\pcPage{Brevet blanc 06}{Mesurer sous l’eau et préserver l’audition}
\phantomsection\label{bb-6}
\pdfbookmark[1]{06 · Mesurer sous l’eau et préserver l’audition}{bm-bb-6}
\noindent\textbf{Volet PC · 30 min · /10}\hfill Partie 1/2 · 4 points scientifiques
\par\vspace{2mm}Les essais et démarches sont pris en compte. Répondre sur une copie séparée.\par
\begin{pcEncadre}{definition}{Le sonar d’un bateau}Un sonar émet un signal ultrasonore de fréquence 40 kHz. Dans l’eau étudiée, sa vitesse est 1 500 m/s. L’écho provenant du fond revient après 0,080 s. Le fond est supposé horizontal et le trajet vertical.\end{pcEncadre}
\pcFigureEx{figures/rev-bb6-sonar-ex.pdf}{Le délai de 0,080 s comprend les deux trajets dans l’eau.}
\par\addvspace{3mm}\noindent\textbf{Question 1 · 1 point}\quad Convertir la fréquence en Hz et la comparer à la limite audible usuelle de 20 000 Hz.\par
\par\addvspace{3mm}\noindent\textbf{Question 2 · 1 point}\quad Quelle distance totale le signal a-t-il parcourue ?\par
\par\addvspace{3mm}\noindent\textbf{Question 3 · 2 points}\quad Calculer la profondeur et justifier le facteur 2.\par
\vfill\pcRepere{cor-bb-6}{Corrigé, barème et remédiation}
\pcPage{Brevet blanc 06}{Mesurer sous l’eau et préserver l’audition}
\noindent\textbf{Volet PC · 30 min · /10}\hfill Partie 2/2 · 5 points scientifiques
\par\vspace{2mm}Les essais et démarches sont pris en compte. Répondre sur une copie séparée.\par
\begin{pcEncadre}{definition}{Alimenter et interpréter}L’appareil fonctionne en courant continu sous 12 V avec une intensité moyenne supposée constante de 0,50 A. Il est utilisé 20 min. On donne $P=UI$, $E=P\Delta t$. Un son peut être intense même si sa fréquence est élevée ou faible.\end{pcEncadre}
\par\addvspace{3mm}\noindent\textbf{Question 4 · 1 point}\quad Calculer la puissance reçue.\par
\par\addvspace{3mm}\noindent\textbf{Question 5 · 1 point}\quad Convertir 20 min en secondes.\par
\par\addvspace{3mm}\noindent\textbf{Question 6 · 1 point}\quad Calculer l’énergie reçue en J.\par
\par\addvspace{3mm}\noindent\textbf{Question 7 · 2 points}\quad En quatre à six lignes, expliquer pourquoi la fréquence seule ne mesure pas le risque auditif et citer deux mesures de prévention adaptées à une écoute sonore.\par
\vfill\pcRepere{cor-bb-6}{Corrigé, barème et remédiation}
{\small 1 point de langue sur l’ensemble de la copie : voir la grille d’autoévaluation.}
\pcPage{Brevet blanc 07}{Recycler et identifier des matériaux}
\phantomsection\label{bb-7}
\pdfbookmark[1]{07 · Recycler et identifier des matériaux}{bm-bb-7}
\noindent\textbf{Volet PC · 30 min · /10}\hfill Partie 1/2 · 4 points scientifiques
\par\vspace{2mm}Les essais et démarches sont pris en compte. Répondre sur une copie séparée.\par
\begin{pcEncadre}{definition}{Un morceau de métal}Un bloc compact de masse 54,0 g déplace 20,0 mL d’eau quand il est complètement immergé sans bulle. Références à température donnée : aluminium 2,70 ; fer 7,8 ; cuivre 8,9 g/cm³. Un noyau d’aluminium décrit ici contient 13 protons et 14 neutrons. Les lectures sont 30,0 mL avant et 50,0 mL après immersion.\end{pcEncadre}
\pcFigureEx{figures/rev-bb7-enonce.pdf}{Les lectures de 30,0 et 50,0 mL correspondent au déplacement de 20,0 mL annoncé.}
\par\addvspace{3mm}\noindent\textbf{Question 1 · 1 point}\quad Calculer la masse volumique du bloc.\par
\par\addvspace{3mm}\noindent\textbf{Question 2 · 1 point}\quad Quelle référence est la plus compatible avec le résultat ?\par
\par\addvspace{3mm}\noindent\textbf{Question 3 · 2 points}\quad Donner Z, A et le nombre d’électrons de l’atome neutre décrit.\par
\vfill\pcRepere{cor-bb-7}{Corrigé, barème et remédiation}
\pcPage{Brevet blanc 07}{Recycler et identifier des matériaux}
\noindent\textbf{Volet PC · 30 min · /10}\hfill Partie 2/2 · 5 points scientifiques
\par\vspace{2mm}Les essais et démarches sont pris en compte. Répondre sur une copie séparée.\par
\begin{pcEncadre}{definition}{Refondre ou transformer}Un atelier pédagogique étudie un modèle de recyclage. Pour un lot, il reçoit 5,0 MJ d’énergie et en transfère 3,5 MJ utilement à la matière. Le reste est transféré vers l’environnement. La fusion est supposée ne former aucune nouvelle espèce. Les données sont fictives.\end{pcEncadre}
\par\addvspace{3mm}\noindent\textbf{Question 4 · 1 point}\quad La fusion décrite est-elle une transformation physique ou chimique ?\par
\par\addvspace{3mm}\noindent\textbf{Question 5 · 1 point}\quad Calculer l’énergie transférée vers l’environnement.\par
\par\addvspace{3mm}\noindent\textbf{Question 6 · 1 point}\quad Une masse de 54,0 g fond sans perte ni réaction : quelle masse liquide obtient-on ?\par
\par\addvspace{3mm}\noindent\textbf{Question 7 · 2 points}\quad En quatre à six lignes, expliquer pourquoi « l’énergie perdue » n’est pas détruite et proposer une limite au slogan « recycler ne coûte aucune ressource ».\par
\vfill\pcRepere{cor-bb-7}{Corrigé, barème et remédiation}
{\small 1 point de langue sur l’ensemble de la copie : voir la grille d’autoévaluation.}
\pcPage{Brevet blanc 08}{Produire et économiser l’électricité}
\phantomsection\label{bb-8}
\pdfbookmark[1]{08 · Produire et économiser l’électricité}{bm-bb-8}
\noindent\textbf{Volet PC · 30 min · /10}\hfill Partie 1/2 · 4 points scientifiques
\par\vspace{2mm}Les essais et démarches sont pris en compte. Répondre sur une copie séparée.\par
\begin{pcEncadre}{definition}{Une journée de production}Le graphique donne la puissance électrique d’un panneau sur une journée modèle. On l’approxime à 0 W avant 10 h, 200 W de 10 h à 14 h, puis 0 W après 14 h. L’énergie est stockée avec un dispositif dont les pertes sont négligées pour ce premier calcul.\end{pcEncadre}
\pcFigureEx{figures/rev-bb8-solaire-ex.pdf}{Repérer les instants de changement de puissance.}
\par\addvspace{3mm}\noindent\textbf{Question 1 · 1 point}\quad Identifier la conversion principale du panneau.\par
\par\addvspace{3mm}\noindent\textbf{Question 2 · 1 point}\quad Lire la durée de production dans le modèle.\par
\par\addvspace{3mm}\noindent\textbf{Question 3 · 2 points}\quad Calculer l’énergie fournie en Wh puis en kWh.\par
\vfill\pcRepere{cor-bb-8}{Corrigé, barème et remédiation}
\pcPage{Brevet blanc 08}{Produire et économiser l’électricité}
\noindent\textbf{Volet PC · 30 min · /10}\hfill Partie 2/2 · 5 points scientifiques
\par\vspace{2mm}Les essais et démarches sont pris en compte. Répondre sur une copie séparée.\par
\begin{pcEncadre}{definition}{Réduire une consommation}Ancienne lampe : 40 W ; nouvelle lampe : 8 W pour le même service d’éclairage étudié. Utilisation : 5 h par jour pendant 100 jours. Prix fictif : 0,25 € par kWh. On donne $E=P\Delta t$.\end{pcEncadre}
\par\addvspace{3mm}\noindent\textbf{Question 4 · 1 point}\quad Calculer la différence de puissance.\par
\par\addvspace{3mm}\noindent\textbf{Question 5 · 1 point}\quad Calculer l’énergie économisée en kWh sur les 100 jours.\par
\par\addvspace{3mm}\noindent\textbf{Question 6 · 1 point}\quad Calculer l’économie monétaire associée.\par
\par\addvspace{3mm}\noindent\textbf{Question 7 · 2 points}\quad En quatre à six lignes, expliquer pourquoi une ressource renouvelable n’assure pas une production constante et proposer un autre critère pour comparer les deux lampes.\par
\vfill\pcRepere{cor-bb-8}{Corrigé, barème et remédiation}
{\small 1 point de langue sur l’ensemble de la copie : voir la grille d’autoévaluation.}
\pcPage{Brevet blanc 09}{Observer un système planétaire}
\phantomsection\label{bb-9}
\pdfbookmark[1]{09 · Observer un système planétaire}{bm-bb-9}
\noindent\textbf{Volet PC · 30 min · /10}\hfill Partie 1/2 · 4 points scientifiques
\par\vspace{2mm}Les essais et démarches sont pris en compte. Répondre sur une copie séparée.\par
\begin{pcEncadre}{definition}{Une lumière ancienne}Une étoile est située à 4 années-lumière. On fournit $1\unite{al}=9{,}46\times10^{12}\unite{km}$ et $c=3{,}00\times10^5\unite{km/s}$. Dans un schéma, une planète reçoit la lumière de l’étoile et en diffuse une partie vers un observateur.\end{pcEncadre}
\pcFigureEx{figures/bb09-etoile.pdf}{}
\par\addvspace{3mm}\noindent\textbf{Question 1 · 1 point}\quad L’année-lumière est-elle une durée ou une distance ?\par
\par\addvspace{3mm}\noindent\textbf{Question 2 · 1 point}\quad Convertir 4 al en km.\par
\par\addvspace{3mm}\noindent\textbf{Question 3 · 2 points}\quad Distinguer étoile et planète comme sources lumineuses et préciser l’ancienneté approximative de la lumière observée de l’étoile.\par
\vfill\pcRepere{cor-bb-9}{Corrigé, barème et remédiation}
\pcPage{Brevet blanc 09}{Observer un système planétaire}
\noindent\textbf{Volet PC · 30 min · /10}\hfill Partie 2/2 · 5 points scientifiques
\par\vspace{2mm}Les essais et démarches sont pris en compte. Répondre sur une copie séparée.\par
\begin{pcEncadre}{definition}{Attraction entre deux corps}Deux corps sont modélisés par $m_A=1{,}0\times10^{20}\unite{kg}$, $m_B=2{,}0\times10^{10}\unite{kg}$ et $d=1{,}0\times10^6\unite m$. On donne $F=Gm_A m_B/d^2$, avec $G=6{,}67\times10^{-11}$ en SI.\end{pcEncadre}
\pcFigureEx{figures/rev-bb9-grav-ex.pdf}{Distance entre centres ; aucune taille n’est à l’échelle.}
\par\addvspace{3mm}\noindent\textbf{Question 4 · 1 point}\quad La force gravitationnelle est-elle attractive ou répulsive ?\par
\par\addvspace{3mm}\noindent\textbf{Question 5 · 1 point}\quad Calculer sa valeur.\par
\par\addvspace{3mm}\noindent\textbf{Question 6 · 1 point}\quad Que devient cette valeur si la distance triple à masses inchangées ?\par
\par\addvspace{3mm}\noindent\textbf{Question 7 · 2 points}\quad En quatre à six lignes, expliquer pourquoi les deux corps subissent des forces de même valeur sans que ces forces se compensent sur un seul corps, puis citer une limite du dessin.\par
\vfill\pcRepere{cor-bb-9}{Corrigé, barème et remédiation}
{\small 1 point de langue sur l’ensemble de la copie : voir la grille d’autoévaluation.}
\pcPage{Brevet blanc 10}{Mener une enquête scientifique}
\phantomsection\label{bb-10}
\pdfbookmark[1]{10 · Mener une enquête scientifique}{bm-bb-10}
\noindent\textbf{Volet PC · 30 min · /10}\hfill Partie 1/2 · 4 points scientifiques
\par\vspace{2mm}Les essais et démarches sont pris en compte. Répondre sur une copie séparée.\par
\begin{pcEncadre}{definition}{Lire une caractéristique}Le graphique U(I) provient d’une résistance scolaire. Un point correspond à I = 0,020 A et U = 4,0 V. On fournit $U=RI$ et $P=UI$ en courant continu. La puissance maximale admissible est de 0,25 W.\end{pcEncadre}
\pcFigureEx{figures/bb10-ohm.pdf}{}
\par\addvspace{3mm}\noindent\textbf{Question 1 · 1 point}\quad Quel aspect du graphe est compatible avec un conducteur ohmique ?\par
\par\addvspace{3mm}\noindent\textbf{Question 2 · 1 point}\quad Calculer la résistance avec le point fourni.\par
\par\addvspace{3mm}\noindent\textbf{Question 3 · 2 points}\quad Calculer la puissance au point donné et comparer à la limite.\par
\vfill\pcRepere{cor-bb-10}{Corrigé, barème et remédiation}
\pcPage{Brevet blanc 10}{Mener une enquête scientifique}
\noindent\textbf{Volet PC · 30 min · /10}\hfill Partie 2/2 · 5 points scientifiques
\par\vspace{2mm}Les essais et démarches sont pris en compte. Répondre sur une copie séparée.\par
\begin{pcEncadre}{definition}{Distinguer les observations}Un professeur présente deux transformations : A, eau liquide qui bout ; B, carbonate réagissant avec un acide et dégageant un gaz. Le gaz B trouble l’eau de chaux. Dans les deux cas, des bulles apparaissent. Une pesée de B en système ouvert passe de 120,0 g à 118,9 g sans projection ni autre échange.\end{pcEncadre}
\pcFigureEx{figures/rev-bb10-transfos-ex.pdf}{Deux transformations différentes peuvent produire des bulles.}
\par\addvspace{3mm}\noindent\textbf{Question 4 · 1 point}\quad Nommer le changement d’état dans A.\par
\par\addvspace{3mm}\noindent\textbf{Question 5 · 1 point}\quad Identifier le gaz mis en évidence dans B.\par
\par\addvspace{3mm}\noindent\textbf{Question 6 · 1 point}\quad Calculer la masse de gaz qui a quitté B avec les hypothèses données.\par
\par\addvspace{3mm}\noindent\textbf{Question 7 · 2 points}\quad En quatre à six lignes, expliquer pourquoi des bulles ne prouvent pas seules une réaction et proposer un contrôle permettant de renforcer l’interprétation de B.\par
\vfill\pcRepere{cor-bb-10}{Corrigé, barème et remédiation}
{\small 1 point de langue sur l’ensemble de la copie : voir la grille d’autoévaluation.}
\global\otCorrigestrue\global\otSoupletrue\gdef\otAccent{otGreenDark}\pcAnnexe{Corrigés des brevets blancs}{Des points de méthode, des réponses et des critères de réussite}\phantomsection\label{cor-brevets}\pdfbookmark[0]{Corrigés des brevets blancs}{bm-cor-brevets}
\begin{pcEncadre}{method}{Se corriger sans masquer ses erreurs}Attribuer les points à partir des critères. Conserver une démarche pertinente même si une erreur de calcul se propage ; éviter de retirer deux fois les mêmes points pour une même erreur. La proposition de correction décrit une répartition pédagogique, pas un barème d’examen.\end{pcEncadre}
\pcSousTitre{Attribuer le point de langue}
Pour chacune des quatre dimensions, attribuer 0 si la compréhension est très gênée ou le critère absent, 0,125 si la maîtrise reste fragile et 0,25 si elle est satisfaisante. Une copie courte correcte n’a pas à devenir artificiellement longue. Prendre en compte les réponses de toute la copie, particulièrement la dernière question rédigée.
\pcSousTitre{Reprendre après le bilan}
Choisir une erreur de connaissance, de modèle ou de conversion. Relire la partie de cours indiquée, refaire un exercice puis reprendre la question du sujet blanc quelques jours plus tard. Une note totale renseigne moins bien qu’une liste précise des points à retravailler.
\pcPage{Corrigé du sujet 01}{Un vélo en mouvement}
\phantomsection\label{cor-bb-1}
\pcRepere{bb-1}{Revenir à l’énoncé}
\par\noindent\begin{minipage}{\linewidth}\pcSousTitre{Question 1 · 1 point}
Rectiligne uniforme : la position croît linéairement avec le temps. \textbf{Barème :} 0,5 point pour rectiligne ; 0,5 pour uniforme et sa lecture graphique.\par
\smallskip\noindent\textbf{Démarche.}\enspace Comparer les déplacements sur des durées égales, dans le référentiel de la route annoncé.\par
\end{minipage}\par
\par\noindent\begin{minipage}{\linewidth}\pcSousTitre{Question 2 · 1 point}
$v=d/\Delta t=200/40=5\unite{m/s}$. \textbf{Barème :} 0,5 pour relation et substitution ; 0,5 pour valeur et unité.\par
\smallskip\noindent\textbf{Démarche.}\enspace Lire les extrémités : $\Delta x=200-0\unite m$ et $\Delta t=40-0\unite s$ ; le quotient donne des m/s.\par
\pcFigureEx{figures/rev-bb1-corrige.pdf}{La pente représente la distance parcourue par unité de temps.}
\end{minipage}\par
\par\noindent\begin{minipage}{\linewidth}\pcSousTitre{Question 3 · 2 points}
$E_c=0{,}5\times80\times5^2=1\,000\unite J$. Si v double à masse fixée, $E_c$ est multipliée par quatre : $4\,000\unite J$. \textbf{Barème :} 1 point pour le premier calcul ; 1 pour facteur et seconde valeur.\par
\smallskip\noindent\textbf{Démarche.}\enspace Employer kg et m/s ; mettre 5 au carré avant le produit. Doubler v multiplie son carré par quatre.\par
\end{minipage}\par
\par\noindent\begin{minipage}{\linewidth}
\begin{pcEncadre}{method}{Ma grille de bilan après les sept questions}Partie 1 : /4 ; partie 2 : /5 ; langue : /1. Total : /10. Après la correction du sujet entier, noter une réussite et une priorité de travail.\end{pcEncadre}
{\small\textbf{Pour reprendre le cours : }\hyperref[chap-3]{chapitre 3 (p.~\pageref*{chap-3})} ; \hyperref[chap-6]{chapitre 6 (p.~\pageref*{chap-6})} ; \hyperref[chap-8]{chapitre 8 (p.~\pageref*{chap-8})} ; \hyperref[chap-11]{chapitre 11 (p.~\pageref*{chap-11})} ; \hyperref[chap-12]{chapitre 12 (p.~\pageref*{chap-12})}.\par}\end{minipage}\par
\pcPage{Corrigé du sujet 01}{Deux solutions énergétiques}
\pcRepere{bb-1}{Revenir à l’énoncé}
\par\noindent\begin{minipage}{\linewidth}\pcSousTitre{Question 4 · 1 point}
Le dioxygène, $\ce{O2}$. \textbf{Barème :} 1 point pour le nom.\par
\smallskip\noindent\textbf{Démarche.}\enspace Le méthane est le combustible ; le dioxygène joue ici le rôle de comburant.\par
\end{minipage}\par
\par\noindent\begin{minipage}{\linewidth}\pcSousTitre{Question 5 · 1 point}
Quatre O de chaque côté : quatre à gauche dans $2\,\ce{O2}$ ; deux dans CO2 et deux dans $2\,\ce{H2O}$ à droite. \textbf{Barème :} 0,5 point par côté.\par
\smallskip\noindent\textbf{Démarche.}\enspace Le coefficient multiplie toute la molécule : $2\,\ce{O2}$ contient $2\times2$ atomes O ; $2\,\ce{H2O}$ en contient $2\times1$.\par
\pcFigureEx{figures/cor-reaction-methane.pdf}{Quatre atomes d’oxygène de chaque côté.}
\end{minipage}\par
\par\noindent\begin{minipage}{\linewidth}\pcSousTitre{Question 6 · 1 point}
$20\times50=1\,000\unite g=1{,}0\unite{kg}$. \textbf{Barème :} 0,5 pour produit ; 0,5 pour conversion et unité.\par
\smallskip\noindent\textbf{Démarche.}\enspace kWh × g/kWh donne g ; diviser ensuite par 1 000 pour obtenir kg.\par
\end{minipage}\par
\par\noindent\begin{minipage}{\linewidth}\pcSousTitre{Question 7 · 2 points}
A est associée à 1,0 kg de CO2 contre 4,0 kg pour B. Dans les conditions et le périmètre des données, A émet quatre fois moins de CO2. Le résultat ne mesure pas tous les impacts, par exemple matériaux, autres polluants ou biodiversité. \textbf{Barème scientifique :} 1 point pour comparaison chiffrée ; 1 pour périmètre et limite pertinente. La langue est appréciée séparément sur l’ensemble du sujet.\par
\smallskip\noindent\textbf{Démarche.}\enspace Commencer par 1,0 kg contre 4,0 kg, puis préciser le périmètre et un impact que l’indicateur ne mesure pas.\par
\end{minipage}\par
\pcPage{Corrigé du sujet 02}{Masse volumique et mélange}
\phantomsection\label{cor-bb-2}
\pcRepere{bb-2}{Revenir à l’énoncé}
\par\noindent\begin{minipage}{\linewidth}\pcSousTitre{Question 1 · 1 point}
$m=82{,}0-32{,}0=50{,}0\unite g$. \textbf{Barème :} 0,5 pour soustraction ; 0,5 pour résultat et unité.\par
\smallskip\noindent\textbf{Démarche.}\enspace La balance a pesé le récipient et le liquide : soustraire la masse du récipient.\par
\end{minipage}\par
\par\noindent\begin{minipage}{\linewidth}\pcSousTitre{Question 2 · 1 point}
$\rho=m/V=50{,}0/50{,}0=1{,}00\unite{g/mL}$. \textbf{Barème :} 0,5 pour quotient ; 0,5 pour valeur et unité.\par
\smallskip\noindent\textbf{Démarche.}\enspace Former $\rho=m/V$ avec la masse du liquide seul ; g et mL donnent g/mL.\par
\pcFigureEx{figures/rev-bb2-masse-cor.pdf}{Soustraire la masse du récipient avant le quotient.}
\end{minipage}\par
\par\noindent\begin{minipage}{\linewidth}\pcSousTitre{Question 3 · 2 points}
Non. Les ions dissous montrent qu’il ne s’agit pas d’un corps pur constitué uniquement d’eau ; une masse volumique voisine de l’eau ne prouve pas la pureté. La potabilité demande des critères chimiques et microbiologiques supplémentaires. \textbf{Barème :} 1 point pour pureté argumentée ; 1 pour potabilité distincte.\par
\smallskip\noindent\textbf{Démarche.}\enspace Les espèces dissoutes concernent la pureté ; les critères sanitaires concernent la potabilité. Donner les deux arguments.\par
\end{minipage}\par
\par\noindent\begin{minipage}{\linewidth}
\begin{pcEncadre}{method}{Ma grille de bilan après les sept questions}Partie 1 : /4 ; partie 2 : /5 ; langue : /1. Total : /10. Après la correction du sujet entier, noter une réussite et une priorité de travail.\end{pcEncadre}
{\small\textbf{Pour reprendre le cours : }\hyperref[chap-1]{chapitre 1 (p.~\pageref*{chap-1})} ; \hyperref[chap-2]{chapitre 2 (p.~\pageref*{chap-2})} ; \hyperref[chap-4]{chapitre 4 (p.~\pageref*{chap-4})}.\par}\end{minipage}\par
\pcPage{Corrigé du sujet 02}{Les résultats d’analyse}
\pcRepere{bb-2}{Revenir à l’énoncé}
\par\noindent\begin{minipage}{\linewidth}\pcSousTitre{Question 4 · 1 point}
Acide, car pH 6 est inférieur à 7 vers 25 °C. \textbf{Barème :} 0,5 pour classement ; 0,5 pour critère.\par
\smallskip\noindent\textbf{Démarche.}\enspace Comparer 6 à 7 ; la proximité avec 7 ne suffit pas à conclure à la neutralité.\par
\end{minipage}\par
\par\noindent\begin{minipage}{\linewidth}\pcSousTitre{Question 5 · 1 point}
L’ion chlorure $\ce{Cl-}$. \textbf{Barème :} 1 point pour nom ou formule non ambiguë.\par
\smallskip\noindent\textbf{Démarche.}\enspace Croiser réactif et couleur du précipité avec la banque de tests ; la couleur initiale du liquide ne remplace pas le test.\par
\end{minipage}\par
\par\noindent\begin{minipage}{\linewidth}\pcSousTitre{Question 6 · 1 point}
18 électrons : Z = 17 protons et un électron de plus pour la charge −1. \textbf{Barème :} 0,5 pour 18 ; 0,5 pour justification.\par
\smallskip\noindent\textbf{Démarche.}\enspace Cl\textsuperscript{-} conserve le même noyau que Cl et gagne un électron : $17+1=18$.\par
\end{minipage}\par
\par\noindent\begin{minipage}{\linewidth}\pcSousTitre{Question 7 · 2 points}
Utiliser des prélèvements séparés, des réactifs adaptés et des témoins. Respecter les protections et éviter de contaminer les flacons. Le test cuivre négatif indique l’absence de précipité détectable dans les conditions ; une quantité inférieure à la sensibilité du test reste possible. \textbf{Barème scientifique :} 1 point pour protocole et témoin ; 1 pour limite du négatif.\par
\smallskip\noindent\textbf{Démarche.}\enspace Prévoir prélèvements séparés, témoin et conditions comparables ; conclure en tenant compte de la limite de détection.\par
\pcFigureEx{figures/rev-bb2-tests-cor.pdf}{Un test négatif ne prouve pas une absence absolue.}
\end{minipage}\par
\pcPage{Corrigé du sujet 03}{Le circuit d’éclairage}
\phantomsection\label{cor-bb-3}
\pcRepere{bb-3}{Revenir à l’énoncé}
\par\noindent\begin{minipage}{\linewidth}\pcSousTitre{Question 1 · 1 point}
12 V pour chaque lampe, car elles sont reliées aux mêmes deux nœuds. \textbf{Barème :} 0,5 pour valeurs ; 0,5 pour loi.\par
\smallskip\noindent\textbf{Démarche.}\enspace Repérer les deux nœuds communs ; chaque lampe est directement reliée à la source par des interrupteurs fermés supposés idéaux.\par
\end{minipage}\par
\par\noindent\begin{minipage}{\linewidth}\pcSousTitre{Question 2 · 1 point}
$I_G=0{,}30+0{,}20=0{,}50\unite A$. \textbf{Barème :} 0,5 pour loi ; 0,5 pour calcul.\par
\smallskip\noindent\textbf{Démarche.}\enspace Les courants sont déjà en ampères : appliquer l’additivité des intensités au nœud.\par
\end{minipage}\par
\par\noindent\begin{minipage}{\linewidth}\pcSousTitre{Question 3 · 2 points}
Voltmètre en dérivation aux bornes de la première ; ampèremètre en série dans la seconde branche, après coupure avant modification. \textbf{Barème :} 1 point par branchement correct.\par
\smallskip\noindent\textbf{Démarche.}\enspace Insérer A dans la branche de L2 hors tension ; relier les deux bornes de V aux deux bornes de L1. A ne se branche pas seul aux bornes de la source.\par
\pcFigureEx{figures/rev-bb3-cor.pdf}{V aux bornes de L1 ; A dans la branche de L2. Chaque interrupteur reste en série avec sa lampe.}
\end{minipage}\par
\par\noindent\begin{minipage}{\linewidth}
\begin{pcEncadre}{method}{Ma grille de bilan après les sept questions}Partie 1 : /4 ; partie 2 : /5 ; langue : /1. Total : /10. Après la correction du sujet entier, noter une réussite et une priorité de travail.\end{pcEncadre}
{\small\textbf{Pour reprendre le cours : }\hyperref[chap-12]{chapitre 12 (p.~\pageref*{chap-12})} ; \hyperref[chap-13]{chapitre 13 (p.~\pageref*{chap-13})} ; \hyperref[chap-14]{chapitre 14 (p.~\pageref*{chap-14})} ; \hyperref[chap-15]{chapitre 15 (p.~\pageref*{chap-15})}.\par}\end{minipage}\par
\pcPage{Corrigé du sujet 03}{Une réserve limitée}
\pcRepere{bb-3}{Revenir à l’énoncé}
\par\noindent\begin{minipage}{\linewidth}\pcSousTitre{Question 4 · 1 point}
$E=6\times4=24\unite{Wh}$. \textbf{Barème :} 0,5 pour relation ; 0,5 pour valeur et unité.\par
\smallskip\noindent\textbf{Démarche.}\enspace Avec W et h, $E=P\Delta t$ donne des Wh.\par
\end{minipage}\par
\par\noindent\begin{minipage}{\linewidth}\pcSousTitre{Question 5 · 1 point}
$24\times3\,600=86\,400\unite J$. \textbf{Barème :} 0,5 pour conversion ; 0,5 pour résultat.\par
\smallskip\noindent\textbf{Démarche.}\enspace Multiplier par $3\,600\unite{J/Wh}$ ; on retrouve le résultat avec $6\times(4\times3\,600)$.\par
\end{minipage}\par
\par\noindent\begin{minipage}{\linewidth}\pcSousTitre{Question 6 · 1 point}
$\Delta t=36/6=6\unite h$. \textbf{Barème :} 0,5 pour division ; 0,5 pour durée et hypothèse.\par
\smallskip\noindent\textbf{Démarche.}\enspace Isoler $\Delta t=E/P$ ; Wh divisés par W donnent h. Contrôler $6\times6=36\unite{Wh}$.\par
\pcFigureEx{figures/rev-bb3-batterie.pdf}{Comparer réserve et consommation dans la même unité.}
\end{minipage}\par
\par\noindent\begin{minipage}{\linewidth}\pcSousTitre{Question 7 · 2 points}
La durée théorique suppose la réserve intégralement utilisable et les seules lampes comme récepteurs. Pertes et autres appareils peuvent réduire l’autonomie. Éteindre une lampe inutile diminue la puissance totale, donc l’énergie reçue pour une même durée. \textbf{Barème scientifique :} 1 point pour une limite liée au modèle ; 1 pour action reliée à E = PΔt.\par
\smallskip\noindent\textbf{Démarche.}\enspace Associer une limite à une hypothèse du modèle, puis une économie à $E=P\Delta t$ ; réduire la durée est aussi pertinent.\par
\end{minipage}\par
\pcPage{Corrigé du sujet 04}{Peser un équipement}
\phantomsection\label{cor-bb-4}
\pcRepere{bb-4}{Revenir à l’énoncé}
\par\noindent\begin{minipage}{\linewidth}\pcSousTitre{Question 1 · 1 point}
12 kg : le transport ne change pas la masse dans ce contexte. \textbf{Barème :} 1 point pour valeur et distinction.\par
\smallskip\noindent\textbf{Démarche.}\enspace La masse caractérise l’équipement ; g et le poids changent avec le lieu.\par
\end{minipage}\par
\par\noindent\begin{minipage}{\linewidth}\pcSousTitre{Question 2 · 1 point}
$P_L=12\times1{,}6=19{,}2\unite N$. \textbf{Barème :} 0,5 pour relation ; 0,5 pour calcul.\par
\smallskip\noindent\textbf{Démarche.}\enspace Choisir $g_L$ ; kg × N/kg donne N. La valeur terrestre ne convient pas ici.\par
\end{minipage}\par
\par\noindent\begin{minipage}{\linewidth}\pcSousTitre{Question 3 · 2 points}
Poids vertical local vers le centre lunaire et action du sol en sens opposé, de valeurs 19,2 N dans le modèle à deux forces. \textbf{Barème :} 1 point pour poids décrit ; 1 pour support et compensation.\par
\smallskip\noindent\textbf{Démarche.}\enspace Nommer les auteurs des actions sur l’équipement : Lune et sol. Dans ce modèle d’équilibre à deux forces, valeurs égales et sens opposés.\par
\pcFigureEx{figures/rev-bb4-forces.pdf}{Deux forces de 19,2 N appliquées à l’équipement ; dessin sans échelle.}
\end{minipage}\par
\par\noindent\begin{minipage}{\linewidth}
\begin{pcEncadre}{method}{Ma grille de bilan après les sept questions}Partie 1 : /4 ; partie 2 : /5 ; langue : /1. Total : /10. Après la correction du sujet entier, noter une réussite et une priorité de travail.\end{pcEncadre}
{\small\textbf{Pour reprendre le cours : }\hyperref[chap-9]{chapitre 9 (p.~\pageref*{chap-9})} ; \hyperref[chap-10]{chapitre 10 (p.~\pageref*{chap-10})} ; \hyperref[chap-16]{chapitre 16 (p.~\pageref*{chap-16})}.\par}\end{minipage}\par
\pcPage{Corrigé du sujet 04}{Communiquer avec la Terre}
\pcRepere{bb-4}{Revenir à l’énoncé}
\par\noindent\begin{minipage}{\linewidth}\pcSousTitre{Question 4 · 1 point}
Électromagnétique ; il peut se propager dans le vide. \textbf{Barème :} 1 point pour nature.\par
\smallskip\noindent\textbf{Démarche.}\enspace Le signal radio est un rayonnement électromagnétique ; un son demande un milieu matériel.\par
\end{minipage}\par
\par\noindent\begin{minipage}{\linewidth}\pcSousTitre{Question 5 · 1 point}
$\Delta t=d/c=1{,}28\unite s$. \textbf{Barème :} 0,5 pour relation ; 0,5 pour résultat.\par
\smallskip\noindent\textbf{Démarche.}\enspace Écrire $\Delta t=d/c$ ; les puissances $10^8$ se simplifient et l’unité obtenue est s.\par
\end{minipage}\par
\par\noindent\begin{minipage}{\linewidth}\pcSousTitre{Question 6 · 1 point}
$2\times1{,}28=2{,}56\unite s$. \textbf{Barème :} 0,5 pour aller-retour ; 0,5 pour durée.\par
\smallskip\noindent\textbf{Démarche.}\enspace Additionner deux trajets de 1,28 s ; ne pas diviser de nouveau une durée calculée pour un aller simple.\par
\pcFigureEx{figures/rev-bb4-radio-cor.pdf}{Chaque trajet mesure 384 000 km ; durée totale : 2,56 s.}
\end{minipage}\par
\par\noindent\begin{minipage}{\linewidth}\pcSousTitre{Question 7 · 2 points}
En orbite, la station et ses occupants restent attirés par l’astre et sont en chute commune. L’impesanteur ne signifie donc pas une force gravitationnelle nulle. Avec les mêmes masses et une distance centrale doublée, le dénominateur est multiplié par quatre ; F est divisée par quatre. \textbf{Barème scientifique :} 1 point pour explication orbitale ; 1 pour facteur justifié.\par
\smallskip\noindent\textbf{Démarche.}\enspace Expliquer la chute commune ; remplacer ensuite d par 2d : $(2d)^2=4d^2$, donc F est divisée par quatre.\par
\end{minipage}\par
\pcPage{Corrigé du sujet 05}{Comparer des protections}
\phantomsection\label{cor-bb-5}
\pcRepere{bb-5}{Revenir à l’énoncé}
\par\noindent\begin{minipage}{\linewidth}\pcSousTitre{Question 1 · 1 point}
Chimique : des espèces de la rouille se forment. \textbf{Barème :} 0,5 pour nature ; 0,5 pour formation d’espèces.\par
\smallskip\noindent\textbf{Démarche.}\enspace La formation de nouvelles espèces distingue cette transformation chimique d’un changement d’état.\par
\end{minipage}\par
\par\noindent\begin{minipage}{\linewidth}\pcSousTitre{Question 2 · 1 point}
Le rôle du dioxygène disponible, l’eau étant présente dans les deux cas. \textbf{Barème :} 1 point pour facteur.\par
\smallskip\noindent\textbf{Démarche.}\enspace B et D contiennent de l’eau ; le facteur comparé est la disponibilité du dioxygène.\par
\end{minipage}\par
\par\noindent\begin{minipage}{\linewidth}\pcSousTitre{Question 3 · 2 points}
Non : « rouille » est qualitatif et ne mesure pas une vitesse. Comparer une quantité de métal corrodé sur des durées identiques, avec répétitions, permettrait d’estimer un effet du sel. \textbf{Barème :} 1 point pour limite ; 1 pour mesure quantitative et conditions.\par
\smallskip\noindent\textbf{Démarche.}\enspace Une observation oui/non ne quantifie pas « deux fois plus vite ». Mesurer une quantité sur une durée avec répétitions et contrôles.\par
\pcFigureEx{figures/rev-cor-corrosion.pdf}{La présence de rouille ne donne pas à elle seule une vitesse.}
\end{minipage}\par
\par\noindent\begin{minipage}{\linewidth}
\begin{pcEncadre}{method}{Ma grille de bilan après les sept questions}Partie 1 : /4 ; partie 2 : /5 ; langue : /1. Total : /10. Après la correction du sujet entier, noter une réussite et une priorité de travail.\end{pcEncadre}
{\small\textbf{Pour reprendre le cours : }\hyperref[chap-2]{chapitre 2 (p.~\pageref*{chap-2})} ; \hyperref[chap-3]{chapitre 3 (p.~\pageref*{chap-3})} ; \hyperref[chap-4]{chapitre 4 (p.~\pageref*{chap-4})} ; \hyperref[chap-5]{chapitre 5 (p.~\pageref*{chap-5})}.\par}\end{minipage}\par
\pcPage{Corrigé du sujet 05}{Le fer dans une solution acide}
\pcRepere{bb-5}{Revenir à l’énoncé}
\par\noindent\begin{minipage}{\linewidth}\pcSousTitre{Question 4 · 1 point}
Dihydrogène, H2. \textbf{Barème :} 1 point.\par
\smallskip\noindent\textbf{Démarche.}\enspace Lire les produits à droite de l’équation : H₂ désigne le dihydrogène.\par
\end{minipage}\par
\par\noindent\begin{minipage}{\linewidth}\pcSousTitre{Question 5 · 1 point}
$26-2=24$ électrons. \textbf{Barème :} 0,5 pour valeur ; 0,5 pour perte de deux électrons.\par
\smallskip\noindent\textbf{Démarche.}\enspace La charge $2+$ signifie deux électrons en moins que l’atome neutre ; Z ne change pas.\par
\pcFigureEx{figures/rev-cor-atome-fe2.pdf}{26 protons et 24 électrons : charge +2.}
\end{minipage}\par
\par\noindent\begin{minipage}{\linewidth}\pcSousTitre{Question 6 · 1 point}
Non : pH 4 est encore acide vers 25 °C. \textbf{Barème :} 0,5 pour réponse ; 0,5 pour critère.\par
\smallskip\noindent\textbf{Démarche.}\enspace Le pH augmente, mais sa valeur finale 4 reste inférieure à 7.\par
\end{minipage}\par
\par\noindent\begin{minipage}{\linewidth}\pcSousTitre{Question 7 · 2 points}
Le gaz H2 et les ions Fe2+ sont les produits de l’équation. La consommation d’ions H+ est cohérente avec une hausse du pH. La corrosion habituelle dans l’air humide met notamment en jeu eau et dioxygène et forme d’autres composés ; l’équation acide-fer ne la décrit pas à elle seule. \textbf{Barème scientifique :} 1 point pour lien observations-modèle ; 1 pour distinction de situations.\par
\smallskip\noindent\textbf{Démarche.}\enspace Relier les observations aux espèces de l’équation, puis distinguer cette réaction de la formation de rouille dans l’air humide.\par
\pcFigureEx{figures/rev-bb5-reaction.pdf}{Les ions H+ sont consommés ; le pH augmente sans atteindre la neutralité.}
\end{minipage}\par
\pcPage{Corrigé du sujet 06}{Le sonar d’un bateau}
\phantomsection\label{cor-bb-6}
\pcRepere{bb-6}{Revenir à l’énoncé}
\par\noindent\begin{minipage}{\linewidth}\pcSousTitre{Question 1 · 1 point}
$40\unite{kHz}=40\,000\unite{Hz}$, au-dessus de la limite usuelle : ultrason. \textbf{Barème :} 0,5 pour conversion ; 0,5 pour comparaison.\par
\smallskip\noindent\textbf{Démarche.}\enspace Kilo signifie mille. Le classement ultrasonore porte sur la fréquence, sans renseigner un niveau en décibels.\par
\end{minipage}\par
\par\noindent\begin{minipage}{\linewidth}\pcSousTitre{Question 2 · 1 point}
$d_{total}=1\,500\times0{,}080=120\unite m$. \textbf{Barème :} 0,5 pour relation ; 0,5 pour calcul.\par
\smallskip\noindent\textbf{Démarche.}\enspace Utiliser la vitesse dans l’eau et la durée totale de 0,080 s pour la longueur totale parcourue.\par
\end{minipage}\par
\par\noindent\begin{minipage}{\linewidth}\pcSousTitre{Question 3 · 2 points}
Le signal parcourt l’aller et le retour, donc la profondeur vaut $120/2=60\unite m$. \textbf{Barème :} 1 point pour trajet double ; 1 pour profondeur et unité.\par
\smallskip\noindent\textbf{Démarche.}\enspace Écrire $2h=v\Delta t$ ; les deux trajets ont la même longueur dans le modèle annoncé.\par
\pcFigureEx{figures/rev-bb6-sonar-cor.pdf}{120 m parcourus au total ; profondeur : 60 m.}
\end{minipage}\par
\par\noindent\begin{minipage}{\linewidth}
\begin{pcEncadre}{method}{Ma grille de bilan après les sept questions}Partie 1 : /4 ; partie 2 : /5 ; langue : /1. Total : /10. Après la correction du sujet entier, noter une réussite et une priorité de travail.\end{pcEncadre}
{\small\textbf{Pour reprendre le cours : }\hyperref[chap-14]{chapitre 14 (p.~\pageref*{chap-14})} ; \hyperref[chap-15]{chapitre 15 (p.~\pageref*{chap-15})} ; \hyperref[chap-17]{chapitre 17 (p.~\pageref*{chap-17})} ; \hyperref[chap-18]{chapitre 18 (p.~\pageref*{chap-18})}.\par}\end{minipage}\par
\pcPage{Corrigé du sujet 06}{Alimenter et interpréter}
\pcRepere{bb-6}{Revenir à l’énoncé}
\par\noindent\begin{minipage}{\linewidth}\pcSousTitre{Question 4 · 1 point}
$P=12\times0{,}50=6{,}0\unite W$. \textbf{Barème :} 0,5 pour relation ; 0,5 pour calcul.\par
\smallskip\noindent\textbf{Démarche.}\enspace En courant continu, $P=UI$ : V × A donne W.\par
\end{minipage}\par
\par\noindent\begin{minipage}{\linewidth}\pcSousTitre{Question 5 · 1 point}
$20\times60=1\,200\unite s$. \textbf{Barème :} 1 point pour conversion.\par
\smallskip\noindent\textbf{Démarche.}\enspace Multiplier par 60 : chaque minute contient 60 secondes.\par
\end{minipage}\par
\par\noindent\begin{minipage}{\linewidth}\pcSousTitre{Question 6 · 1 point}
$E=6\times1\,200=7\,200\unite J$. \textbf{Barème :} 0,5 pour relation ; 0,5 pour résultat.\par
\smallskip\noindent\textbf{Démarche.}\enspace W × s donne J ; ne pas conserver 20 minutes comme si c’étaient 20 secondes.\par
\pcFigureEx{figures/rev-bb6-energie.pdf}{Convertir la durée avant le produit avec la puissance.}
\end{minipage}\par
\par\noindent\begin{minipage}{\linewidth}\pcSousTitre{Question 7 · 2 points}
La fréquence caractérise surtout la hauteur et ne donne pas le niveau sonore ni la durée d’exposition. Le risque dépend notamment de ces deux facteurs. Réduire le niveau, s’éloigner de la source et faire des pauses sont des mesures adaptées ; une protection appropriée peut être nécessaire. \textbf{Barème scientifique :} 1 point pour distinction fréquence/niveau/durée ; 0,5 par mesure cohérente, jusqu’à 1 point.\par
\smallskip\noindent\textbf{Démarche.}\enspace Distinguer hauteur, niveau et durée ; justifier les mesures par la réduction du niveau reçu ou du temps d’exposition.\par
\end{minipage}\par
\pcPage{Corrigé du sujet 07}{Un morceau de métal}
\phantomsection\label{cor-bb-7}
\pcRepere{bb-7}{Revenir à l’énoncé}
\par\noindent\begin{minipage}{\linewidth}\pcSousTitre{Question 1 · 1 point}
$\rho=54{,}0/20{,}0=2{,}70\unite{g/cm^3}$. \textbf{Barème :} 0,5 pour quotient ; 0,5 pour unité et valeur.\par
\smallskip\noindent\textbf{Démarche.}\enspace $V=50-30=20\unite{mL}=20\unite{cm^3}$. Diviser 54 g par ce volume pour obtenir 2,70 g/cm³.\par
\pcFigureEx{figures/rev-bb7-corrige.pdf}{50 − 30 = 20 cm³ ; 54 / 20 = 2,70 g/cm³.}
\end{minipage}\par
\par\noindent\begin{minipage}{\linewidth}\pcSousTitre{Question 2 · 1 point}
L’aluminium ; cette seule propriété ne prouve pas la pureté. \textbf{Barème :} 1 point pour identification compatible.\par
\smallskip\noindent\textbf{Démarche.}\enspace Comparer dans la même unité et à température comparable ; formuler une compatibilité, pas une preuve absolue d’identité.\par
\end{minipage}\par
\par\noindent\begin{minipage}{\linewidth}\pcSousTitre{Question 3 · 2 points}
$Z=13$, $A=27$ et 13 électrons. \textbf{Barème :} 0,5 pour Z ; 0,5 pour A ; 1 pour nombre d’électrons et neutralité.\par
\smallskip\noindent\textbf{Démarche.}\enspace Z compte les protons, A les protons et neutrons. Dans l’atome neutre demandé, le nombre d’électrons égale Z.\par
\pcFigureEx{figures/rev-bb7-atome.pdf}{Inventaire symbolique des constituants de l’atome neutre décrit ; aucune échelle de taille.}
\end{minipage}\par
\par\noindent\begin{minipage}{\linewidth}
\begin{pcEncadre}{method}{Ma grille de bilan après les sept questions}Partie 1 : /4 ; partie 2 : /5 ; langue : /1. Total : /10. Après la correction du sujet entier, noter une réussite et une priorité de travail.\end{pcEncadre}
{\small\textbf{Pour reprendre le cours : }\hyperref[chap-1]{chapitre 1 (p.~\pageref*{chap-1})} ; \hyperref[chap-2]{chapitre 2 (p.~\pageref*{chap-2})} ; \hyperref[chap-3]{chapitre 3 (p.~\pageref*{chap-3})} ; \hyperref[chap-12]{chapitre 12 (p.~\pageref*{chap-12})}.\par}\end{minipage}\par
\pcPage{Corrigé du sujet 07}{Refondre ou transformer}
\pcRepere{bb-7}{Revenir à l’énoncé}
\par\noindent\begin{minipage}{\linewidth}\pcSousTitre{Question 4 · 1 point}
Physique : changement d’état sans nouvelle espèce. \textbf{Barème :} 0,5 pour nature ; 0,5 pour justification.\par
\smallskip\noindent\textbf{Démarche.}\enspace Le document exclut toute nouvelle espèce : la fusion est physique dans ce modèle.\par
\end{minipage}\par
\par\noindent\begin{minipage}{\linewidth}\pcSousTitre{Question 5 · 1 point}
$5{,}0-3{,}5=1{,}5\unite{MJ}$. \textbf{Barème :} 0,5 pour bilan ; 0,5 pour résultat.\par
\smallskip\noindent\textbf{Démarche.}\enspace Sans stockage ni autre sortie, soustraire l’énergie utile de l’entrée. Vérifier $3{,}5+1{,}5=5{,}0\unite{MJ}$.\par
\pcFigureEx{figures/rev-bb7-bilan.pdf}{5,0 MJ = 3,5 MJ utiles + 1,5 MJ vers le milieu.}
\end{minipage}\par
\par\noindent\begin{minipage}{\linewidth}\pcSousTitre{Question 6 · 1 point}
54,0 g, par conservation de la masse dans les conditions données. \textbf{Barème :} 1 point.\par
\smallskip\noindent\textbf{Démarche.}\enspace Un changement d’état sans perte ni réaction conserve la masse.\par
\end{minipage}\par
\par\noindent\begin{minipage}{\linewidth}\pcSousTitre{Question 7 · 2 points}
L’énergie transférée vers l’environnement existe encore, notamment sous forme thermique, mais n’est plus utile au procédé. Le recyclage demande aussi énergie, collecte et tri. Il peut réduire certains prélèvements sans supprimer tous les besoins ni toutes les pertes de matière. \textbf{Barème scientifique :} 1 point pour conservation avec périmètre ; 1 pour besoin de ressources expliqué.\par
\smallskip\noindent\textbf{Démarche.}\enspace Préciser où l’énergie est transférée, puis citer un besoin concret du recyclage. Un transfert thermique ne disparaît pas.\par
\end{minipage}\par
\pcPage{Corrigé du sujet 08}{Une journée de production}
\phantomsection\label{cor-bb-8}
\pcRepere{bb-8}{Revenir à l’énoncé}
\par\noindent\begin{minipage}{\linewidth}\pcSousTitre{Question 1 · 1 point}
Transfert reçu par rayonnement lumineux, conversion en énergie électrique, avec d’autres sorties notamment thermiques dans la réalité. \textbf{Barème :} 1 point pour entrée et sortie principales.\par
\smallskip\noindent\textbf{Démarche.}\enspace Distinguer apport par rayonnement et sortie électrique ; les autres transferts réels restent dans le bilan.\par
\end{minipage}\par
\par\noindent\begin{minipage}{\linewidth}\pcSousTitre{Question 2 · 1 point}
4 h. \textbf{Barème :} 1 point pour lecture de 10 h à 14 h.\par
\smallskip\noindent\textbf{Démarche.}\enspace Lire la largeur du palier : $14-10=4\unite h$ ; 14 h est l’heure de fin, pas la durée.\par
\end{minipage}\par
\par\noindent\begin{minipage}{\linewidth}\pcSousTitre{Question 3 · 2 points}
$E=200\times4=800\unite{Wh}=0{,}800\unite{kWh}$. \textbf{Barème :} 1 point pour calcul Wh ; 1 pour conversion kWh.\par
\smallskip\noindent\textbf{Démarche.}\enspace La puissance constante permet $E=P\Delta t$. Diviser les Wh par 1 000 pour obtenir les kWh.\par
\pcFigureEx{figures/rev-bb8-solaire-cor.pdf}{L’aire du rectangle donne l’énergie : W × h = Wh.}
\end{minipage}\par
\par\noindent\begin{minipage}{\linewidth}
\begin{pcEncadre}{method}{Ma grille de bilan après les sept questions}Partie 1 : /4 ; partie 2 : /5 ; langue : /1. Total : /10. Après la correction du sujet entier, noter une réussite et une priorité de travail.\end{pcEncadre}
{\small\textbf{Pour reprendre le cours : }\hyperref[chap-12]{chapitre 12 (p.~\pageref*{chap-12})} ; \hyperref[chap-13]{chapitre 13 (p.~\pageref*{chap-13})} ; \hyperref[chap-14]{chapitre 14 (p.~\pageref*{chap-14})} ; \hyperref[chap-15]{chapitre 15 (p.~\pageref*{chap-15})}.\par}\end{minipage}\par
\pcPage{Corrigé du sujet 08}{Réduire une consommation}
\pcRepere{bb-8}{Revenir à l’énoncé}
\par\noindent\begin{minipage}{\linewidth}\pcSousTitre{Question 4 · 1 point}
$40-8=32\unite W$. \textbf{Barème :} 1 point.\par
\smallskip\noindent\textbf{Démarche.}\enspace Soustraire les puissances pour le même service : $\Delta P=40-8=32\unite W$.\par
\end{minipage}\par
\par\noindent\begin{minipage}{\linewidth}\pcSousTitre{Question 5 · 1 point}
$E=0{,}032\times5\times100=16\unite{kWh}$. \textbf{Barème :} 0,5 pour unités et durée ; 0,5 pour résultat.\par
\smallskip\noindent\textbf{Démarche.}\enspace La durée vaut $5\times100=500\unite h$ ; convertir 32 W en 0,032 kW avant le produit en kWh.\par
\end{minipage}\par
\par\noindent\begin{minipage}{\linewidth}\pcSousTitre{Question 6 · 1 point}
$16\times0{,}25=4{,}00$ €. \textbf{Barème :} 0,5 pour produit ; 0,5 pour coût.\par
\smallskip\noindent\textbf{Démarche.}\enspace kWh × euros/kWh donne des euros ; ce coût concerne le seul poste fourni.\par
\pcFigureEx{figures/rev-bb8-lampes.pdf}{Chaque lampe fonctionne 500 h ; économie : 16 kWh.}
\end{minipage}\par
\par\noindent\begin{minipage}{\linewidth}\pcSousTitre{Question 7 · 2 points}
La ressource solaire se renouvelle, mais la puissance dépend de l’éclairement, de la météo et de l’heure. Un stockage ou une autre production peut être nécessaire pour certains usages. Durée de vie, qualité de l’éclairage ou impacts de fabrication sont d’autres critères pertinents. \textbf{Barème scientifique :} 1 point pour variabilité expliquée ; 1 pour autre critère cohérent.\par
\smallskip\noindent\textbf{Démarche.}\enspace Relier éclairement et puissance disponible, puis ajouter un critère comme durée de vie ou qualité de lumière.\par
\end{minipage}\par
\pcPage{Corrigé du sujet 09}{Une lumière ancienne}
\phantomsection\label{cor-bb-9}
\pcRepere{bb-9}{Revenir à l’énoncé}
\par\noindent\begin{minipage}{\linewidth}\pcSousTitre{Question 1 · 1 point}
Une distance. \textbf{Barème :} 1 point.\par
\smallskip\noindent\textbf{Démarche.}\enspace L’année-lumière est la distance parcourue par la lumière en un an ; elle reste une unité de distance.\par
\end{minipage}\par
\par\noindent\begin{minipage}{\linewidth}\pcSousTitre{Question 2 · 1 point}
$4\times9{,}46\times10^{12}=3{,}784\times10^{13}\unite{km}$. \textbf{Barème :} 0,5 pour produit ; 0,5 pour résultat et unité.\par
\smallskip\noindent\textbf{Démarche.}\enspace Multiplier la distance d’une année-lumière par quatre ; conserver puissance de dix et unité km.\par
\end{minipage}\par
\par\noindent\begin{minipage}{\linewidth}\pcSousTitre{Question 3 · 2 points}
L’étoile est une source primaire ; la planète est un objet diffusant dans le modèle. La lumière de l’étoile reçue a été émise environ quatre ans auparavant. \textbf{Barème :} 1 point pour distinction ; 1 pour délai.\par
\smallskip\noindent\textbf{Démarche.}\enspace La planète diffuse la lumière reçue. Les quatre ans de retard concernent la distance étoile-observateur fournie.\par
\pcFigureEx{figures/rev-bb9-lumiere.pdf}{La distance en années-lumière indique le retard en années.}
\end{minipage}\par
\par\noindent\begin{minipage}{\linewidth}
\begin{pcEncadre}{method}{Ma grille de bilan après les sept questions}Partie 1 : /4 ; partie 2 : /5 ; langue : /1. Total : /10. Après la correction du sujet entier, noter une réussite et une priorité de travail.\end{pcEncadre}
{\small\textbf{Pour reprendre le cours : }\hyperref[chap-7]{chapitre 7 (p.~\pageref*{chap-7})} ; \hyperref[chap-10]{chapitre 10 (p.~\pageref*{chap-10})} ; \hyperref[chap-16]{chapitre 16 (p.~\pageref*{chap-16})}.\par}\end{minipage}\par
\pcPage{Corrigé du sujet 09}{Attraction entre deux corps}
\pcRepere{bb-9}{Revenir à l’énoncé}
\par\noindent\begin{minipage}{\linewidth}\pcSousTitre{Question 4 · 1 point}
Attractive. \textbf{Barème :} 1 point.\par
\smallskip\noindent\textbf{Démarche.}\enspace La force sur chaque corps est dirigée vers l’autre corps : l’interaction est attractive.\par
\end{minipage}\par
\par\noindent\begin{minipage}{\linewidth}\pcSousTitre{Question 5 · 1 point}
$F=1{,}334\times10^8\unite N$, soit environ $1{,}3\times10^8\unite N$. \textbf{Barème :} 0,5 pour substitution avec d² ; 0,5 pour calcul et unité.\par
\smallskip\noindent\textbf{Démarche.}\enspace Les données sont en SI ; mettre la distance au carré avec sa puissance de dix : $(10^6)^2=10^{12}$.\par
\end{minipage}\par
\par\noindent\begin{minipage}{\linewidth}\pcSousTitre{Question 6 · 1 point}
F est divisée par neuf, soit environ $1{,}48\times10^7\unite N$. \textbf{Barème :} 0,5 pour facteur ; 0,5 pour valeur.\par
\smallskip\noindent\textbf{Démarche.}\enspace Remplacer d par 3d donne F/9 ; utiliser la valeur non arrondie avant de diviser.\par
\end{minipage}\par
\par\noindent\begin{minipage}{\linewidth}\pcSousTitre{Question 7 · 2 points}
Les deux forces appartiennent à la même interaction mais s’exercent sur des corps différents. Elles ne sont donc pas deux forces opposées appliquées au même système. Le dessin n’est pas à l’échelle et le modèle assimile les corps à des masses ponctuelles ou sphériques. \textbf{Barème scientifique :} 1 point pour systèmes distincts ; 1 pour limite annoncée.\par
\smallskip\noindent\textbf{Démarche.}\enspace Identifier le corps qui reçoit chaque force ; une compensation demanderait des forces appliquées au même système.\par
\pcFigureEx{figures/rev-bb9-grav-cor.pdf}{Chaque corps subit 1,334 × 10⁸ N à la distance initiale ; les deux forces s’appliquent à deux systèmes différents.}
\end{minipage}\par
\pcPage{Corrigé du sujet 10}{Lire une caractéristique}
\phantomsection\label{cor-bb-10}
\pcRepere{bb-10}{Revenir à l’énoncé}
\par\noindent\begin{minipage}{\linewidth}\pcSousTitre{Question 1 · 1 point}
Une droite passant par l’origine, à la dispersion des mesures près. \textbf{Barème :} 0,5 pour droite ; 0,5 pour origine.\par
\smallskip\noindent\textbf{Démarche.}\enspace Il faut une droite et le passage par l’origine ; une droite décalée ne suffit pas pour une proportionnalité.\par
\end{minipage}\par
\par\noindent\begin{minipage}{\linewidth}\pcSousTitre{Question 2 · 1 point}
$R=4{,}0/0{,}020=200\unite\Omega$. \textbf{Barème :} 0,5 pour quotient ; 0,5 pour résultat.\par
\smallskip\noindent\textbf{Démarche.}\enspace Choisir un point hors origine, avec I en A ; le quotient U/I donne des ohms.\par
\end{minipage}\par
\par\noindent\begin{minipage}{\linewidth}\pcSousTitre{Question 3 · 2 points}
$P=4{,}0\times0{,}020=0{,}080\unite W$, inférieure à 0,25 W. \textbf{Barème :} 1 point pour calcul ; 1 pour comparaison justifiée.\par
\smallskip\noindent\textbf{Démarche.}\enspace Calculer P au même point de fonctionnement ; comparer 0,080 W à 0,25 W dans la même unité.\par
\pcFigureEx{figures/rev-bb10-ohm-cor.pdf}{R = 200 ohms ; P = 0,080 W, inférieure à 0,25 W.}
\end{minipage}\par
\par\noindent\begin{minipage}{\linewidth}
\begin{pcEncadre}{method}{Ma grille de bilan après les sept questions}Partie 1 : /4 ; partie 2 : /5 ; langue : /1. Total : /10. Après la correction du sujet entier, noter une réussite et une priorité de travail.\end{pcEncadre}
{\small\textbf{Pour reprendre le cours : }\hyperref[chap-3]{chapitre 3 (p.~\pageref*{chap-3})} ; \hyperref[chap-4]{chapitre 4 (p.~\pageref*{chap-4})} ; \hyperref[chap-8]{chapitre 8 (p.~\pageref*{chap-8})} ; \hyperref[chap-14]{chapitre 14 (p.~\pageref*{chap-14})} ; \hyperref[chap-18]{chapitre 18 (p.~\pageref*{chap-18})}.\par}\end{minipage}\par
\pcPage{Corrigé du sujet 10}{Distinguer les observations}
\pcRepere{bb-10}{Revenir à l’énoncé}
\par\noindent\begin{minipage}{\linewidth}\pcSousTitre{Question 4 · 1 point}
Vaporisation, ici ébullition. \textbf{Barème :} 1 point.\par
\smallskip\noindent\textbf{Démarche.}\enspace Nommer le passage liquide vers gaz ; l’ébullition se produit dans le volume du liquide.\par
\end{minipage}\par
\par\noindent\begin{minipage}{\linewidth}\pcSousTitre{Question 5 · 1 point}
Dioxyde de carbone. \textbf{Barème :} 1 point.\par
\smallskip\noindent\textbf{Démarche.}\enspace La conclusion vient du test à l’eau de chaux, pas des seules bulles.\par
\end{minipage}\par
\par\noindent\begin{minipage}{\linewidth}\pcSousTitre{Question 6 · 1 point}
$120{,}0-118{,}9=1{,}1\unite g$. \textbf{Barème :} 0,5 pour bilan ; 0,5 pour masse et unité.\par
\smallskip\noindent\textbf{Démarche.}\enspace La masse sortie vaut masse initiale moins masse finale dans les hypothèses du système ouvert.\par
\end{minipage}\par
\par\noindent\begin{minipage}{\linewidth}\pcSousTitre{Question 7 · 2 points}
Les bulles peuvent provenir d’un changement d’état, d’un dégazage ou d’une réaction. Dans B, le test identifie CO2 ; il faut aussi comparer les espèces initiales et finales. Un témoin sans mise en contact des réactifs et un contrôle de l’eau de chaux renforcent l’attribution. \textbf{Barème scientifique :} 1 point pour mécanismes distincts ; 1 pour contrôle pertinent.\par
\smallskip\noindent\textbf{Démarche.}\enspace Comparer les espèces avant et après ; proposer un témoin sans gaz ajouté et des conditions de test comparables.\par
\pcFigureEx{figures/rev-bb10-transfos-cor.pdf}{Le test au CO2 renforce l’interprétation chimique de B ; les bulles seules ne suffisent pas.}
\end{minipage}\par
\global\otCorrigesfalse\global\otSouplefalse\gdef\otAccent{otNavy}\pcAnnexe{Cinq sujets format annales}{Physique-chimie · 30 minutes · Sujets originaux 11 à 15}\phantomsection\label{annales-originales}\pdfbookmark[0]{Cinq sujets au format des annales}{bm-annales-originales}
\begin{pcEncadre}{method}{Des sujets complets pour composer}Chaque sujet contient trois parties indépendantes, neuf questions et tous les documents nécessaires. Durée conseillée : 30 minutes. Le total est de 9 points scientifiques, complétés par 1 point de langue selon la grille du livre. Les contextes, les mesures, les questions et les figures ont été créés pour ces entraînements. Ce ne sont pas des annales officielles reproduites.\end{pcEncadre}
\pcSousTitre{Un parcours en trois temps}
Prévoir environ 8 minutes par partie, puis 6 minutes pour lire et relire l’ensemble. Une réponse justifiée compte davantage qu’une valeur isolée. Les relations utiles et les hypothèses du modèle figurent dans les documents. Les schémas sans échelle ne doivent pas être mesurés à la règle.
\pcSousTitre{Après la composition}
Pour chaque question, le corrigé distingue la réponse, la démarche et les critères du barème. Garder les étapes justes même si une erreur de calcul se propage. Les points d’une même erreur ne sont pas retirés deux fois. Le point de langue porte sur l’ensemble des réponses, sans imposer un texte inutilement long.
\pcRepere{ann-11}{Sujet 11 · Un monte-charge pour les archives}
\pcRepere{ann-12}{Sujet 12 · Comprendre un détecteur optique de fumée}
\pcRepere{ann-13}{Sujet 13 · Une station météo sur la côte}
\pcRepere{ann-14}{Sujet 14 · Transporter un colis réfrigéré}
\pcRepere{ann-15}{Sujet 15 · Une maquette de dirigeable}
\pcPage{Sujet 11 · Énoncé}{Un monte-charge pour les archives}
\phantomsection\label{ann-11}
\pdfbookmark[1]{11 · Un monte-charge pour les archives}{bm-ann-11}
\textbf{PC · 30 min · /10}\hfill Partie 1/3 · 3 points scientifiques
\begin{pcEncadre}{definition}{Document · Suivre la cabine}Une cabine et son chargement ont une masse totale de 50 kg. La cabine monte verticalement ; son altitude est relevée toutes les secondes. Pendant la portion étudiée, les frottements sont négligés. On modélise la cabine chargée par un point G soumis à son poids et à la force du câble. $g=9{,}8\unite{N/kg}$.\end{pcEncadre}
\pcFigureEx{figures/ann11-positions.pdf}{Document original · données pédagogiques. Schéma sans échelle, sauf axes gradués.}
\par\addvspace{3mm}\noindent\begin{minipage}{\linewidth}\textbf{Question 1 · 0,5 point}\quad Décrire le mouvement de la cabine dans le référentiel du bâtiment en justifiant les deux mots employés.\par\end{minipage}\par
\par\addvspace{3mm}\noindent\begin{minipage}{\linewidth}\textbf{Question 2 · 1 point}\quad Calculer la vitesse moyenne entre 0 s et 3 s.\par\end{minipage}\par
\par\addvspace{3mm}\noindent\begin{minipage}{\linewidth}\textbf{Question 3 · 1,5 points}\quad Calculer le poids de l’ensemble. Pour le mouvement étudié, les forces se compensent. Représenter sans échelle, au point G, les deux forces et indiquer leur valeur.\par\end{minipage}\par
\vfill\pcRepere{cor-ann-11}{Corrigé expliqué, barème et bilan}
\pcPage{Sujet 11 · Énoncé}{Un monte-charge pour les archives}
\textbf{PC · 30 min · /10}\hfill Partie 2/3 · 3 points scientifiques
\begin{pcEncadre}{definition}{Document · Prévoir l’autonomie}Le moteur reçoit une tension continue de 24 V et une intensité de 3,0 A, supposées constantes pendant un cycle complet de 6 min. Une batterie fournit 18 Wh utilisables. Aucun autre appareil ne la sollicite dans le modèle. Relations : $P=UI$, $E=P\Delta t$ ; avec W et h, E est en Wh.\end{pcEncadre}
\par\addvspace{3mm}\noindent\begin{minipage}{\linewidth}\textbf{Question 4 · 1 point}\quad Calculer la puissance électrique reçue par le moteur.\par\end{minipage}\par
\par\addvspace{3mm}\noindent\begin{minipage}{\linewidth}\textbf{Question 5 · 1 point}\quad Calculer l’énergie électrique nécessaire à un cycle en Wh.\par\end{minipage}\par
\par\addvspace{3mm}\noindent\begin{minipage}{\linewidth}\textbf{Question 6 · 1 point}\quad La batterie permet-elle trois cycles complets ? Justifier, puis donner le nombre maximal de cycles dans ce modèle.\par\end{minipage}\par
\vfill\pcRepere{cor-ann-11}{Corrigé expliqué, barème et bilan}
\pcPage{Sujet 11 · Énoncé}{Un monte-charge pour les archives}
\textbf{PC · 30 min · /10}\hfill Partie 3/3 · 3 points scientifiques
\begin{pcEncadre}{definition}{Document · Identifier un liquide de maintenance}Le laboratoire pèse un récipient vide : 28,5 g. Le même récipient contenant 100,0 mL d’un liquide a une masse de 102,5 g. Références à la même température : A, 0,74 g/mL ; B, 0,90 g/mL ; C, 1,15 g/mL. On donne $\rho=m/V$.\end{pcEncadre}
\pcFigureEx{figures/ann11-pesees.pdf}{Document original · données pédagogiques. Schéma sans échelle, sauf axes gradués.}
\par\addvspace{3mm}\noindent\begin{minipage}{\linewidth}\textbf{Question 7 · 0,5 point}\quad Déterminer la masse du liquide seul.\par\end{minipage}\par
\par\addvspace{3mm}\noindent\begin{minipage}{\linewidth}\textbf{Question 8 · 1 point}\quad Calculer la masse volumique du liquide et indiquer la référence compatible.\par\end{minipage}\par
\par\addvspace{3mm}\noindent\begin{minipage}{\linewidth}\textbf{Question 9 · 1,5 points}\quad En trois à cinq lignes, expliquer pourquoi cette compatibilité ne prouve pas l’identité chimique du liquide. Proposer un contrôle supplémentaire pertinent.\par\end{minipage}\par
\vfill\pcRepere{cor-ann-11}{Corrigé expliqué, barème et bilan}
\pcPage{Sujet 12 · Énoncé}{Comprendre un détecteur optique de fumée}
\phantomsection\label{ann-12}
\pdfbookmark[1]{12 · Comprendre un détecteur optique de fumée}{bm-ann-12}
\textbf{PC · 30 min · /10}\hfill Partie 1/3 · 3 points scientifiques
\begin{pcEncadre}{definition}{Document · Un trajet lumineux modifié}Dans ce modèle de détecteur, une DEL émet de la lumière dans une petite chambre. Le capteur est placé hors du faisceau direct. Sans fumée, il ne reçoit presque pas de lumière. Des particules de fumée peuvent diffuser une partie de la lumière vers lui. Le dessin représente la présence de fumée.\end{pcEncadre}
\pcFigureEx{figures/ann12-optique.pdf}{Document original · données pédagogiques. Schéma sans échelle, sauf axes gradués.}
\par\addvspace{3mm}\noindent\begin{minipage}{\linewidth}\textbf{Question 1 · 0,5 point}\quad Identifier la source primaire de lumière et le récepteur.\par\end{minipage}\par
\par\addvspace{3mm}\noindent\begin{minipage}{\linewidth}\textbf{Question 2 · 1 point}\quad Reproduire le schéma simplifié et ajouter les deux flèches nécessaires pour expliquer l’arrivée de lumière au capteur.\par\end{minipage}\par
\par\addvspace{3mm}\noindent\begin{minipage}{\linewidth}\textbf{Question 3 · 1,5 points}\quad En trois à cinq lignes, expliquer pourquoi le capteur est placé hors du faisceau direct. Pourquoi de la poussière peut-elle aussi perturber ce modèle ?\par\end{minipage}\par
\vfill\pcRepere{cor-ann-12}{Corrigé expliqué, barème et bilan}
\pcPage{Sujet 12 · Énoncé}{Comprendre un détecteur optique de fumée}
\textbf{PC · 30 min · /10}\hfill Partie 2/3 · 3 points scientifiques
\begin{pcEncadre}{definition}{Document · Raisonner sur une combustion}Une combustion complète du propane est modélisée par $\ce{C3H8 + 5 O2 -> 3 CO2 + 4 H2O}$. Le document ne modélise pas la formation de fumée ni les produits d’une combustion incomplète.\end{pcEncadre}
\pcFigureEx{figures/ann12-molecules.pdf}{Document original · données pédagogiques. Schéma sans échelle, sauf axes gradués.}
\par\addvspace{3mm}\noindent\begin{minipage}{\linewidth}\textbf{Question 4 · 0,5 point}\quad Nommer le combustible et le comburant de cette réaction.\par\end{minipage}\par
\par\addvspace{3mm}\noindent\begin{minipage}{\linewidth}\textbf{Question 5 · 1 point}\quad Vérifier la conservation des atomes d’oxygène de part et d’autre de l’équation.\par\end{minipage}\par
\par\addvspace{3mm}\noindent\begin{minipage}{\linewidth}\textbf{Question 6 · 1,5 points}\quad Pendant un essai, le réservoir perd 4,4 g de propane. Un élève en déduit que la masse totale des produits formés vaut 4,4 g. Peut-on valider sa conclusion ? Identifier la masse oubliée et écrire le bilan de masse complet, sans calcul numérique.\par\end{minipage}\par
\vfill\pcRepere{cor-ann-12}{Corrigé expliqué, barème et bilan}
\pcPage{Sujet 12 · Énoncé}{Comprendre un détecteur optique de fumée}
\textbf{PC · 30 min · /10}\hfill Partie 3/3 · 3 points scientifiques
\begin{pcEncadre}{definition}{Document · Recevoir le signal sonore}Un essai en plein air utilise un signal bref. Un microphone est situé à 51 m de l’émetteur. On néglige le vent et les réflexions ; la vitesse du son est $340\unite{m/s}$. Deux enregistrements A et B portent sur une même durée de 4 ms.\end{pcEncadre}
\pcFigureEx{figures/ann12-son.pdf}{Document original · données pédagogiques. Schéma sans échelle, sauf axes gradués.}
\par\addvspace{3mm}\noindent\begin{minipage}{\linewidth}\textbf{Question 7 · 1 point}\quad Calculer la durée de propagation directe entre l’émetteur et le microphone.\par\end{minipage}\par
\par\addvspace{3mm}\noindent\begin{minipage}{\linewidth}\textbf{Question 8 · 1 point}\quad Quel enregistrement correspond au son le plus aigu ? Justifier sans confondre fréquence et amplitude.\par\end{minipage}\par
\par\addvspace{3mm}\noindent\begin{minipage}{\linewidth}\textbf{Question 9 · 1 point}\quad Le son pourrait-il se propager de la même manière dans le vide ? Justifier.\par\end{minipage}\par
\vfill\pcRepere{cor-ann-12}{Corrigé expliqué, barème et bilan}
\pcPage{Sujet 13 · Énoncé}{Une station météo sur la côte}
\phantomsection\label{ann-13}
\pdfbookmark[1]{13 · Une station météo sur la côte}{bm-ann-13}
\textbf{PC · 30 min · /10}\hfill Partie 1/3 · 3 points scientifiques
\begin{pcEncadre}{definition}{Document · Connaître l’air prélevé}Un flacon contient 2,0 L d’air sec. Le modèle utilise 78 \% de diazote, 21 \% de dioxygène et 1 \% d’autres gaz en volume. L’azote possède $Z=7$ ; l’isotope étudié possède $A=14$. La formule du diazote est $\ce{N2}$.\end{pcEncadre}
\pcFigureEx{figures/ann13-air.pdf}{Document original · données pédagogiques. Schéma sans échelle, sauf axes gradués.}
\par\addvspace{3mm}\noindent\begin{minipage}{\linewidth}\textbf{Question 1 · 0,5 point}\quad L’air sec est-il un corps pur ou un mélange ? Justifier.\par\end{minipage}\par
\par\addvspace{3mm}\noindent\begin{minipage}{\linewidth}\textbf{Question 2 · 1 point}\quad Calculer le volume de dioxygène du prélèvement, en litre puis en millilitres.\par\end{minipage}\par
\par\addvspace{3mm}\noindent\begin{minipage}{\linewidth}\textbf{Question 3 · 1,5 points}\quad Donner les nombres de protons, de neutrons et d’électrons d’un atome neutre d’azote de cet isotope. Combien d’atomes contient une molécule de diazote ?\par\end{minipage}\par
\vfill\pcRepere{cor-ann-13}{Corrigé expliqué, barème et bilan}
\pcPage{Sujet 13 · Énoncé}{Une station météo sur la côte}
\textbf{PC · 30 min · /10}\hfill Partie 2/3 · 3 points scientifiques
\begin{pcEncadre}{definition}{Document · Tester un composant résistif}Le graphique présente des mesures sur un conducteur ohmique. On utilise $U=RI$ avec U en V, I en A et R en ohms. Le montage de principe du corrigé sera réalisé uniquement avec le matériel scolaire adapté.\end{pcEncadre}
\pcFigureEx{figures/ann13-ohm.pdf}{Document original · données pédagogiques. Schéma sans échelle, sauf axes gradués.}
\par\addvspace{3mm}\noindent\begin{minipage}{\linewidth}\textbf{Question 4 · 0,5 point}\quad Lire la tension correspondant à une intensité de 30 mA.\par\end{minipage}\par
\par\addvspace{3mm}\noindent\begin{minipage}{\linewidth}\textbf{Question 5 · 1 point}\quad Calculer la résistance à partir de ce point.\par\end{minipage}\par
\par\addvspace{3mm}\noindent\begin{minipage}{\linewidth}\textbf{Question 6 · 1,5 points}\quad Schématiser un montage permettant ces mesures avec générateur, résistance, ampèremètre et voltmètre. Préciser leurs branchements.\par\end{minipage}\par
\vfill\pcRepere{cor-ann-13}{Corrigé expliqué, barème et bilan}
\pcPage{Sujet 13 · Énoncé}{Une station météo sur la côte}
\textbf{PC · 30 min · /10}\hfill Partie 3/3 · 3 points scientifiques
\begin{pcEncadre}{definition}{Document · Transmettre un message}La station envoie un signal radio à un relais distant de 120 km. La vitesse retenue est $3{,}0\times10^8\unite{m/s}$. Le message d’étude comporte quatre symboles avec le code A = 10, B = 01, C = 11, D = 00 ; il est composé de huit intervalles de 2 ms. La durée de préparation n’est pas incluse.\end{pcEncadre}
\pcFigureEx{figures/ann13-code.pdf}{Document original · données pédagogiques. Schéma sans échelle, sauf axes gradués.}
\par\addvspace{3mm}\noindent\begin{minipage}{\linewidth}\textbf{Question 7 · 0,5 point}\quad Décoder les quatre groupes représentés sur le document.\par\end{minipage}\par
\par\addvspace{3mm}\noindent\begin{minipage}{\linewidth}\textbf{Question 8 · 1 point}\quad Calculer la durée totale d’émission du message.\par\end{minipage}\par
\par\addvspace{3mm}\noindent\begin{minipage}{\linewidth}\textbf{Question 9 · 1,5 points}\quad Calculer la durée de propagation du début du message jusqu’au relais. Expliquer pourquoi elle diffère de la durée d’émission.\par\end{minipage}\par
\vfill\pcRepere{cor-ann-13}{Corrigé expliqué, barème et bilan}
\pcPage{Sujet 14 · Énoncé}{Transporter un colis réfrigéré}
\phantomsection\label{ann-14}
\pdfbookmark[1]{14 · Transporter un colis réfrigéré}{bm-ann-14}
\textbf{PC · 30 min · /10}\hfill Partie 1/3 · 3 points scientifiques
\begin{pcEncadre}{definition}{Document · Étudier une réserve froide}Un sachet étanche contient 200 g d’eau pure initialement solide. Sous la pression usuelle, le professeur suit sa température pendant un apport d’énergie. Le sachet reste fermé, sans fuite ; on mesure ici la température de l’eau dans ses différents états.\end{pcEncadre}
\pcFigureEx{figures/ann14-fusion.pdf}{Document original · données pédagogiques. Schéma sans échelle, sauf axes gradués.}
\par\addvspace{3mm}\noindent\begin{minipage}{\linewidth}\textbf{Question 1 · 1 point}\quad Nommer le changement d’état au palier de 0 °C et indiquer les états qui coexistent.\par\end{minipage}\par
\par\addvspace{3mm}\noindent\begin{minipage}{\linewidth}\textbf{Question 2 · 1 point}\quad Lire la durée du palier et prévoir la masse d’eau après fusion complète.\par\end{minipage}\par
\par\addvspace{3mm}\noindent\begin{minipage}{\linewidth}\textbf{Question 3 · 1 point}\quad Un élève affirme : « La température ne change plus, donc l’eau ne reçoit plus d’énergie. » Discuter cette affirmation.\par\end{minipage}\par
\vfill\pcRepere{cor-ann-14}{Corrigé expliqué, barème et bilan}
\pcPage{Sujet 14 · Énoncé}{Transporter un colis réfrigéré}
\textbf{PC · 30 min · /10}\hfill Partie 2/3 · 3 points scientifiques
\begin{pcEncadre}{definition}{Document · Analyser un liquide de conditionnement}Un autre prélèvement de liquide présente un pH de 8,0 vers 25 °C. La fiche de test indique : nitrate d’argent + ions chlorure → précipité blanc. Après ajout du réactif à un prélèvement séparé, un précipité blanc apparaît. Pour le chlore, $Z=17$. Les tests sont réalisés au laboratoire sous la conduite du professeur.\end{pcEncadre}
\pcFigureEx{figures/ann14-test.pdf}{Document original · données pédagogiques. Schéma sans échelle, sauf axes gradués.}
\par\addvspace{3mm}\noindent\begin{minipage}{\linewidth}\textbf{Question 4 · 0,5 point}\quad Classer le liquide selon son pH et justifier.\par\end{minipage}\par
\par\addvspace{3mm}\noindent\begin{minipage}{\linewidth}\textbf{Question 5 · 1 point}\quad Identifier l’ion mis en évidence et donner son nombre d’électrons.\par\end{minipage}\par
\par\addvspace{3mm}\noindent\begin{minipage}{\linewidth}\textbf{Question 6 · 1,5 points}\quad En trois à cinq lignes, expliquer ce que permettent les résultats et ce qu’ils ne prouvent pas sur la composition et l’innocuité du liquide.\par\end{minipage}\par
\vfill\pcRepere{cor-ann-14}{Corrigé expliqué, barème et bilan}
\pcPage{Sujet 14 · Énoncé}{Transporter un colis réfrigéré}
\textbf{PC · 30 min · /10}\hfill Partie 3/3 · 3 points scientifiques
\begin{pcEncadre}{definition}{Document · Limiter la vitesse sur un convoyeur}Le colis de masse 2,0 kg se déplace en ligne droite sur un convoyeur. Une phase à vitesse constante est comparée à deux réglages : 0,50 m/s et 1,0 m/s. On donne $E_c=\frac12mv^2$, avec les unités SI.\end{pcEncadre}
\pcFigureEx{figures/ann14-convoyeur.pdf}{Document original · données pédagogiques. Schéma sans échelle, sauf axes gradués.}
\par\addvspace{3mm}\noindent\begin{minipage}{\linewidth}\textbf{Question 7 · 1 point}\quad Quelle distance parcourt le colis en 8,0 s à 0,50 m/s constants ?\par\end{minipage}\par
\par\addvspace{3mm}\noindent\begin{minipage}{\linewidth}\textbf{Question 8 · 1 point}\quad Calculer l’énergie cinétique du colis pour chacun des deux réglages.\par\end{minipage}\par
\par\addvspace{3mm}\noindent\begin{minipage}{\linewidth}\textbf{Question 9 · 1 point}\quad Le doublement de vitesse double-t-il l’énergie cinétique ? Justifier, puis expliquer l’intérêt de limiter la vitesse.\par\end{minipage}\par
\vfill\pcRepere{cor-ann-14}{Corrigé expliqué, barème et bilan}
\pcPage{Sujet 15 · Énoncé}{Une maquette de dirigeable}
\phantomsection\label{ann-15}
\pdfbookmark[1]{15 · Une maquette de dirigeable}{bm-ann-15}
\textbf{PC · 30 min · /10}\hfill Partie 1/3 · 3 points scientifiques
\begin{pcEncadre}{definition}{Document · Décrire le gaz de l’enveloppe}La maquette contient de l’hélium modélisé par des atomes neutres séparés. L’isotope représenté possède $Z=2$ et $A=4$. Les points du modèle ne sont pas dessinés à la taille réelle des atomes.\end{pcEncadre}
\pcFigureEx{figures/ann15-helium.pdf}{Document original · données pédagogiques. Schéma sans échelle, sauf axes gradués.}
\par\addvspace{3mm}\noindent\begin{minipage}{\linewidth}\textbf{Question 1 · 1 point}\quad Donner les nombres de protons, de neutrons et d’électrons de cet atome.\par\end{minipage}\par
\par\addvspace{3mm}\noindent\begin{minipage}{\linewidth}\textbf{Question 2 · 1 point}\quad Le modèle représente-t-il un gaz constitué de molécules de plusieurs atomes ? Justifier.\par\end{minipage}\par
\par\addvspace{3mm}\noindent\begin{minipage}{\linewidth}\textbf{Question 3 · 1 point}\quad Après compression douce sans fuite, à température redevenue identique, le nombre d’atomes diminue-t-il ? Expliquer la modification du modèle.\par\end{minipage}\par
\vfill\pcRepere{cor-ann-15}{Corrigé expliqué, barème et bilan}
\pcPage{Sujet 15 · Énoncé}{Une maquette de dirigeable}
\textbf{PC · 30 min · /10}\hfill Partie 2/3 · 3 points scientifiques
\begin{pcEncadre}{definition}{Document · Équilibrer la maquette}La masse totale est de 0,30 kg. Le modèle retient le poids vertical vers le bas et une force de l’air verticale vers le haut de 3,4 N. Le dirigeable peut être retenu par un fil vertical ancré au sol. On donne $g=10\unite{N/kg}$.\end{pcEncadre}
\pcFigureEx{figures/ann15-forces-enonce.pdf}{Document original · données pédagogiques. Schéma sans échelle, sauf axes gradués.}
\par\addvspace{3mm}\noindent\begin{minipage}{\linewidth}\textbf{Question 4 · 0,5 point}\quad Calculer le poids total de la maquette.\par\end{minipage}\par
\par\addvspace{3mm}\noindent\begin{minipage}{\linewidth}\textbf{Question 5 · 1 point}\quad Sans fil, les deux forces indiquées se compensent-elles ? Justifier.\par\end{minipage}\par
\par\addvspace{3mm}\noindent\begin{minipage}{\linewidth}\textbf{Question 6 · 1,5 points}\quad Le fil permet de maintenir le dirigeable immobile. Déterminer la valeur et le sens de sa force sur la maquette, puis représenter les trois forces sans échelle.\par\end{minipage}\par
\vfill\pcRepere{cor-ann-15}{Corrigé expliqué, barème et bilan}
\pcPage{Sujet 15 · Énoncé}{Une maquette de dirigeable}
\textbf{PC · 30 min · /10}\hfill Partie 3/3 · 3 points scientifiques
\begin{pcEncadre}{definition}{Document · Alimenter les accessoires}Un générateur adapté fournit une tension constante de 6,0 V à deux branches en dérivation : moteur, 0,20 A ; éclairage, 0,10 A. Les courants restent constants dans le modèle. La batterie fournit 1,8 Wh utilisables. On donne $P=UI$ et $E=P\Delta t$.\end{pcEncadre}
\pcFigureEx{figures/ann15-circuit.pdf}{Document original · données pédagogiques. Schéma sans échelle, sauf axes gradués.}
\par\addvspace{3mm}\noindent\begin{minipage}{\linewidth}\textbf{Question 7 · 1 point}\quad Donner la tension aux bornes de chaque branche et calculer l’intensité fournie par le générateur.\par\end{minipage}\par
\par\addvspace{3mm}\noindent\begin{minipage}{\linewidth}\textbf{Question 8 · 1 point}\quad Calculer la puissance totale et la durée maximale d’alimentation dans ce modèle.\par\end{minipage}\par
\par\addvspace{3mm}\noindent\begin{minipage}{\linewidth}\textbf{Question 9 · 1 point}\quad La branche de l’éclairage devient ouverte. Le moteur est-il nécessairement arrêté ? Expliquer avec le schéma.\par\end{minipage}\par
\vfill\pcRepere{cor-ann-15}{Corrigé expliqué, barème et bilan}
\global\otCorrigestrue\global\otSoupletrue\gdef\otAccent{otGreenDark}\pcAnnexe{Corrigés des sujets 11 à 15}{Comprendre le modèle, les calculs et les critères de réussite}\phantomsection\label{cor-annales-originales}\pdfbookmark[0]{Corrigés des sujets 11 à 15}{bm-cor-annales-originales}
\begin{pcEncadre}{method}{Une grille pour progresser}Chaque sujet : partie 1 /3, partie 2 /3, partie 3 /3 et langue /1, soit /10. Relire d’abord la démarche avant d’attribuer les points. Repérer une erreur de modèle, une erreur de lecture et une erreur d’unité si elles existent, puis choisir un exercice du cours pour reprendre la difficulté.\end{pcEncadre}
\pcPage{Corrigé du sujet 11}{Suivre la cabine}
\phantomsection\label{cor-ann-11}
\pcRepere{ann-11}{Revenir à l’énoncé}
\par\noindent\begin{minipage}{\linewidth}\pcSousTitre{Question 1}Le mouvement est rectiligne, car la trajectoire est verticale et droite, et uniforme : l’altitude augmente de 0,60 m pendant chaque seconde.\par\smallskip\textbf{Démarche.}\enspace Séparer la forme de la trajectoire et l’évolution de la valeur de la vitesse. Les relevés sont faits à durées égales.\par\smallskip{\small\textbf{Barème.}\enspace 0,25 pour rectiligne justifié ; 0,25 pour uniforme justifié.\par}\end{minipage}\par
\par\noindent\begin{minipage}{\linewidth}\pcSousTitre{Question 2}$v=d/\Delta t=1{,}80/3{,}0=0{,}60\unite{m/s}$. La distance parcourue est 1,80 m pendant cette montée.\par\smallskip\textbf{Démarche.}\enspace Utiliser la différence entre les altitudes initiale et finale, puis la durée du même intervalle.\par\smallskip{\small\textbf{Barème.}\enspace 0,5 pour quotient et substitution ; 0,5 pour valeur et unité.\par}\end{minipage}\par
\par\noindent\begin{minipage}{\linewidth}\pcSousTitre{Question 3}$P=mg=50\times9{,}8=490\unite N$. La force du câble vaut aussi 490 N : verticale vers le haut, opposée au poids vertical vers le bas. Elles agissent ici sur le même système.\par\smallskip\textbf{Démarche.}\enspace La compensation impose deux forces de même direction, de même valeur et de sens opposés. Elle est compatible avec une vitesse constante non nulle.\par\smallskip{\small\textbf{Barème.}\enspace 0,5 pour le poids ; 0,5 pour directions et sens ; 0,5 pour égalité des valeurs et représentation.\par}\end{minipage}\par
\pcFigureEx{figures/ann11-forces.pdf}{Modèle de correction · les légendes précisent les directions, les connexions ou les résultats.}
\pcPage{Corrigé du sujet 11}{Prévoir l’autonomie}
\pcRepere{ann-11}{Revenir à l’énoncé}
\par\noindent\begin{minipage}{\linewidth}\pcSousTitre{Question 4}$P=24\times3{,}0=72\unite W$.\par\smallskip\textbf{Démarche.}\enspace Multiplier tension en volts et intensité en ampères. Le watt mesure une puissance.\par\smallskip{\small\textbf{Barème.}\enspace 0,5 pour relation et substitution ; 0,5 pour résultat avec unité.\par}\end{minipage}\par
\par\noindent\begin{minipage}{\linewidth}\pcSousTitre{Question 5}$6\unite{min}=6/60\unite h=0{,}10\unite h$ ; $E=72\times0{,}10=7{,}2\unite{Wh}$.\par\smallskip\textbf{Démarche.}\enspace Convertir les minutes en heures avant de multiplier une puissance en watts pour obtenir des Wh.\par\smallskip{\small\textbf{Barème.}\enspace 0,5 pour conversion ; 0,5 pour produit et unité.\par}\end{minipage}\par
\par\noindent\begin{minipage}{\linewidth}\pcSousTitre{Question 6}Trois cycles demandent $3\times7{,}2=21{,}6\unite{Wh}$, au-delà des 18 Wh disponibles. Deux cycles demandent 14,4 Wh : le maximum est donc deux cycles complets.\par\smallskip\textbf{Démarche.}\enspace Comparer des énergies dans la même unité, puis retenir un nombre entier de cycles. Une fraction de cycle ne constitue pas un cycle complet.\par\smallskip{\small\textbf{Barème.}\enspace 0,5 pour comparaison chiffrée ; 0,5 pour deux cycles justifiés.\par}\end{minipage}\par
\pcFigureEx{figures/ann11-autonomie.pdf}{Modèle de correction · les légendes précisent les directions, les connexions ou les résultats.}
\pcPage{Corrigé du sujet 11}{Identifier un liquide de maintenance}
\pcRepere{ann-11}{Revenir à l’énoncé}
\par\noindent\begin{minipage}{\linewidth}\pcSousTitre{Question 7}$m=102{,}5-28{,}5=74{,}0\unite g$.\par\smallskip\textbf{Démarche.}\enspace La masse totale comprend le récipient et son contenu : retrancher la masse du récipient.\par\smallskip{\small\textbf{Barème.}\enspace 0,25 pour soustraction ; 0,25 pour valeur et unité.\par}\end{minipage}\par
\par\noindent\begin{minipage}{\linewidth}\pcSousTitre{Question 8}$\rho=74{,}0/100{,}0=0{,}740\unite{g/mL}$. La valeur est compatible avec A, aux arrondis des références près.\par\smallskip\textbf{Démarche.}\enspace Conserver les unités g et mL pour comparer directement aux références du document.\par\smallskip{\small\textbf{Barème.}\enspace 0,5 pour quotient et unité ; 0,5 pour comparaison à A.\par}\end{minipage}\par
\par\noindent\begin{minipage}{\linewidth}\pcSousTitre{Question 9}Une même masse volumique peut correspondre à plusieurs substances ou à un mélange. La mesure possède aussi une incertitude. Il faut comparer une autre propriété indépendante, par exemple une température de changement d’état dans un protocole adapté, ou réaliser une analyse chimique appropriée au laboratoire.\par\smallskip\textbf{Démarche.}\enspace Distinguer « compatible avec » et « prouvé ». Le contrôle doit porter sur une propriété différente et rester adapté au produit étudié.\par\smallskip{\small\textbf{Barème.}\enspace 0,5 pour possibilité de substances ou mélanges différents ; 0,5 pour limites de mesure ; 0,5 pour contrôle indépendant pertinent.\par}\end{minipage}\par
\par\noindent\begin{minipage}{\linewidth}{\small\textbf{Bilan :} /3 + /3 + /3 + /1 = /10.\par\textbf{Pour reprendre :} \hyperref[chap-1]{chapitre 1 (p.~\pageref*{chap-1})} ; \hyperref[chap-8]{chapitre 8 (p.~\pageref*{chap-8})} ; \hyperref[chap-9]{chapitre 9 (p.~\pageref*{chap-9})} ; \hyperref[chap-14]{chapitre 14 (p.~\pageref*{chap-14})} ; \hyperref[chap-15]{chapitre 15 (p.~\pageref*{chap-15})}.\par}\end{minipage}\par
\pcPage{Corrigé du sujet 12}{Un trajet lumineux modifié}
\phantomsection\label{cor-ann-12}
\pcRepere{ann-12}{Revenir à l’énoncé}
\par\noindent\begin{minipage}{\linewidth}\pcSousTitre{Question 1}La DEL est la source primaire : elle produit la lumière émise. Le capteur est le récepteur de lumière.\par\smallskip\textbf{Démarche.}\enspace Attribuer les rôles en suivant la lumière, de sa production jusqu’à sa détection.\par\smallskip{\small\textbf{Barème.}\enspace 0,25 pour la DEL ; 0,25 pour le capteur.\par}\end{minipage}\par
\par\noindent\begin{minipage}{\linewidth}\pcSousTitre{Question 2}La première flèche va de la DEL aux particules ; la seconde va des particules vers le capteur. Les particules diffusent une partie de la lumière reçue.\par\smallskip\textbf{Démarche.}\enspace Les flèches indiquent le sens de propagation de la lumière. Le capteur n’envoie pas de rayon pour chercher la fumée.\par\smallskip{\small\textbf{Barème.}\enspace 0,5 par trajet fléché correctement orienté.\par}\end{minipage}\par
\par\noindent\begin{minipage}{\linewidth}\pcSousTitre{Question 3}Hors du faisceau direct, le capteur reçoit peu de lumière en l’absence de particules. Une augmentation de lumière reçue peut donc indiquer une diffusion par des particules. La poussière peut également diffuser de la lumière : cette seule mesure ne donne pas l’identité chimique des particules.\par\smallskip\textbf{Démarche.}\enspace Relier le placement du capteur au signal de référence, puis séparer détection d’une présence et identification de sa nature.\par\smallskip{\small\textbf{Barème.}\enspace 0,5 pour faible signal sans particules ; 0,5 pour diffusion ; 0,5 pour limite liée à la poussière.\par}\end{minipage}\par
\pcFigureEx{figures/ann12-optique-cor.pdf}{Modèle de correction · les légendes précisent les directions, les connexions ou les résultats.}
\pcPage{Corrigé du sujet 12}{Raisonner sur une combustion}
\pcRepere{ann-12}{Revenir à l’énoncé}
\par\noindent\begin{minipage}{\linewidth}\pcSousTitre{Question 4}Le propane est le combustible ; le dioxygène est le comburant.\par\smallskip\textbf{Démarche.}\enspace Repérer la substance qui brûle et le réactif qui permet la combustion.\par\smallskip{\small\textbf{Barème.}\enspace 0,25 pour chacun des deux noms.\par}\end{minipage}\par
\par\noindent\begin{minipage}{\linewidth}\pcSousTitre{Question 5}À gauche, 5 molécules de dioxygène apportent 10 atomes O. À droite, 3 molécules de dioxyde de carbone apportent 6 atomes O et 4 molécules d’eau en apportent 4 : total 10.\par\smallskip\textbf{Démarche.}\enspace Multiplier le nombre devant la formule par le nombre d’atomes O dans une molécule, puis additionner les produits.\par\smallskip{\small\textbf{Barème.}\enspace 0,5 pour le décompte à gauche ; 0,5 pour celui à droite.\par}\end{minipage}\par
\par\noindent\begin{minipage}{\linewidth}\pcSousTitre{Question 6}Non : la réaction consomme aussi du dioxygène prélevé dans l’air. Sa masse doit être ajoutée à celle du propane consommé. Pour la transformation complète du document : $m_{\rm propane}+m_{\rm dioxygène}=m_{\rm dioxyde\ de\ carbone}+m_{\rm eau}$. La seule perte du réservoir ne donne pas la masse totale des produits.\par\smallskip\textbf{Démarche.}\enspace Définir la frontière du système pesé : le réservoir ne contient pas tout le dioxygène consommé. La conservation porte sur tous les réactifs et tous les produits, pas sur la seule masse de combustible.\par\smallskip{\small\textbf{Barème.}\enspace 0,5 pour rejet justifié ; 0,5 pour dioxygène oublié ; 0,5 pour bilan complet.\par}\end{minipage}\par
\pcFigureEx{figures/ann12-bilan.pdf}{Modèle de correction · les légendes précisent les directions, les connexions ou les résultats.}
\pcPage{Corrigé du sujet 12}{Recevoir le signal sonore}
\pcRepere{ann-12}{Revenir à l’énoncé}
\par\noindent\begin{minipage}{\linewidth}\pcSousTitre{Question 7}$\Delta t=d/v=51/340=0{,}15\unite s$.\par\smallskip\textbf{Démarche.}\enspace Il s’agit d’un trajet simple : aucune division par deux n’intervient.\par\smallskip{\small\textbf{Barème.}\enspace 0,5 pour quotient ; 0,5 pour durée et unité.\par}\end{minipage}\par
\par\noindent\begin{minipage}{\linewidth}\pcSousTitre{Question 8}B correspond au son le plus aigu : il présente quatre oscillations pendant les 4 ms, contre deux pour A. Les amplitudes dessinées sont égales.\par\smallskip\textbf{Démarche.}\enspace Comparer le nombre d’oscillations sur une même durée. Une plus grande amplitude ne signifie pas un son plus aigu.\par\smallskip{\small\textbf{Barème.}\enspace 0,5 pour B ; 0,5 pour comparaison à durée égale.\par}\end{minipage}\par
\par\noindent\begin{minipage}{\linewidth}\pcSousTitre{Question 9}Non. Le son est une perturbation mécanique qui nécessite un milieu matériel pour se propager. Le vide ne fournit pas ce milieu.\par\smallskip\textbf{Démarche.}\enspace Distinguer la propagation du son de celle de la lumière, qui peut se propager dans le vide.\par\smallskip{\small\textbf{Barème.}\enspace 0,5 pour réponse ; 0,5 pour milieu matériel nécessaire.\par}\end{minipage}\par
\pcFigureEx{figures/ann12-son-cor.pdf}{Modèle de correction · les légendes précisent les directions, les connexions ou les résultats.}
\par\noindent\begin{minipage}{\linewidth}{\small\textbf{Bilan :} /3 + /3 + /3 + /1 = /10.\par\textbf{Pour reprendre :} \hyperref[chap-3]{chapitre 3 (p.~\pageref*{chap-3})} ; \hyperref[chap-6]{chapitre 6 (p.~\pageref*{chap-6})} ; \hyperref[chap-16]{chapitre 16 (p.~\pageref*{chap-16})} ; \hyperref[chap-17]{chapitre 17 (p.~\pageref*{chap-17})}.\par}\end{minipage}\par
\pcPage{Corrigé du sujet 13}{Connaître l’air prélevé}
\phantomsection\label{cor-ann-13}
\pcRepere{ann-13}{Revenir à l’énoncé}
\par\noindent\begin{minipage}{\linewidth}\pcSousTitre{Question 1}C’est un mélange : il contient plusieurs espèces chimiques, dont le diazote et le dioxygène.\par\smallskip\textbf{Démarche.}\enspace La transparence et l’aspect uniforme ne prouvent pas qu’un gaz est un corps pur.\par\smallskip{\small\textbf{Barème.}\enspace 0,25 pour mélange ; 0,25 pour pluralité des espèces.\par}\end{minipage}\par
\par\noindent\begin{minipage}{\linewidth}\pcSousTitre{Question 2}$V=0{,}21\times2{,}0=0{,}42\unite L=420\unite{mL}$.\par\smallskip\textbf{Démarche.}\enspace Appliquer la proportion au volume total, puis convertir les litres en millilitres.\par\smallskip{\small\textbf{Barème.}\enspace 0,5 pour produit ; 0,5 pour valeurs et conversion.\par}\end{minipage}\par
\par\noindent\begin{minipage}{\linewidth}\pcSousTitre{Question 3}L’atome possède 7 protons, 14 − 7 = 7 neutrons et 7 électrons car il est neutre. Une molécule de diazote contient deux atomes d’azote.\par\smallskip\textbf{Démarche.}\enspace Z compte les protons ; A compte protons et neutrons. La neutralité impose autant d’électrons que de protons. L’indice 2 porte sur les atomes de la molécule.\par\smallskip{\small\textbf{Barème.}\enspace 0,25 pour chacun des trois inventaires ; 0,25 pour neutralité justifiée ; 0,5 pour deux atomes.\par}\end{minipage}\par
\pcFigureEx{figures/ann13-atome.pdf}{Modèle de correction · les légendes précisent les directions, les connexions ou les résultats.}
\pcPage{Corrigé du sujet 13}{Tester un composant résistif}
\pcRepere{ann-13}{Revenir à l’énoncé}
\par\noindent\begin{minipage}{\linewidth}\pcSousTitre{Question 4}$U=6{,}0\unite V$ pour $I=30\unite{mA}$.\par\smallskip\textbf{Démarche.}\enspace Partir de 30 mA sur l’axe horizontal, rejoindre la droite puis lire la tension sur l’axe vertical.\par\smallskip{\small\textbf{Barème.}\enspace 0,5 pour lecture et unité.\par}\end{minipage}\par
\par\noindent\begin{minipage}{\linewidth}\pcSousTitre{Question 5}$30\unite{mA}=0{,}030\unite A$ ; $R=6{,}0/0{,}030=200\unite\Omega$.\par\smallskip\textbf{Démarche.}\enspace Convertir les milliampères en ampères avant d’appliquer la loi d’Ohm.\par\smallskip{\small\textbf{Barème.}\enspace 0,5 pour conversion et quotient ; 0,5 pour résultat et unité.\par}\end{minipage}\par
\par\noindent\begin{minipage}{\linewidth}\pcSousTitre{Question 6}L’ampèremètre se place en série dans la branche de la résistance ; le voltmètre se branche en dérivation entre les deux bornes de la résistance. Les bornes A et V sont du côté des potentiels les plus élevés, les bornes COM de l’autre côté, pour des lectures positives.\par\smallskip\textbf{Démarche.}\enspace Mesurer I impose de traverser l’ampèremètre. Mesurer U impose de comparer deux points du circuit : le voltmètre relie ces deux points.\par\smallskip{\small\textbf{Barème.}\enspace 0,5 pour boucle et résistance ; 0,5 pour ampèremètre en série ; 0,5 pour voltmètre en dérivation.\par}\end{minipage}\par
\pcFigureEx{figures/ann13-mesures.pdf}{Modèle de correction · les légendes précisent les directions, les connexions ou les résultats.}
\pcPage{Corrigé du sujet 13}{Transmettre un message}
\pcRepere{ann-13}{Revenir à l’énoncé}
\par\noindent\begin{minipage}{\linewidth}\pcSousTitre{Question 7}Le message est CADB : 11 → C, 10 → A, 00 → D, 01 → B.\par\smallskip\textbf{Démarche.}\enspace Lire chaque groupe depuis le début annoncé ; ne pas décaler le découpage en groupes de deux.\par\smallskip{\small\textbf{Barème.}\enspace 0,5 pour les quatre symboles dans l’ordre.\par}\end{minipage}\par
\par\noindent\begin{minipage}{\linewidth}\pcSousTitre{Question 8}$8\times2=16\unite{ms}=0{,}016\unite s$.\par\smallskip\textbf{Démarche.}\enspace Compter les huit intervalles, et non seulement les intervalles codés 1. Le zéro occupe lui aussi du temps.\par\smallskip{\small\textbf{Barème.}\enspace 0,5 pour huit intervalles ; 0,5 pour produit et unité.\par}\end{minipage}\par
\par\noindent\begin{minipage}{\linewidth}\pcSousTitre{Question 9}$120\unite{km}=120\,000\unite m$ ; $\Delta t=120\,000/(3{,}0\times10^8)=4{,}0\times10^{-4}\unite s=0{,}40\unite{ms}$. Cette durée est le trajet d’un front du signal, tandis que les 16 ms correspondent à l’émission des huit valeurs successives.\par\smallskip\textbf{Démarche.}\enspace Faire le calcul en mètres et secondes. Un message long peut commencer à arriver avant que son émission soit terminée.\par\smallskip{\small\textbf{Barème.}\enspace 0,5 pour conversion ; 0,5 pour quotient, valeur et unité ; 0,5 pour distinction des durées.\par}\end{minipage}\par
\pcFigureEx{figures/ann13-radio.pdf}{Modèle de correction · les légendes précisent les directions, les connexions ou les résultats.}
\par\noindent\begin{minipage}{\linewidth}{\small\textbf{Bilan :} /3 + /3 + /3 + /1 = /10.\par\textbf{Pour reprendre :} \hyperref[chap-2]{chapitre 2 (p.~\pageref*{chap-2})} ; \hyperref[chap-6]{chapitre 6 (p.~\pageref*{chap-6})} ; \hyperref[chap-13]{chapitre 13 (p.~\pageref*{chap-13})} ; \hyperref[chap-14]{chapitre 14 (p.~\pageref*{chap-14})} ; \hyperref[chap-18]{chapitre 18 (p.~\pageref*{chap-18})}.\par}\end{minipage}\par
\pcPage{Corrigé du sujet 14}{Étudier une réserve froide}
\phantomsection\label{cor-ann-14}
\pcRepere{ann-14}{Revenir à l’énoncé}
\par\noindent\begin{minipage}{\linewidth}\pcSousTitre{Question 1}Il s’agit de la fusion. Pendant ce palier, l’eau solide et l’eau liquide coexistent dans le modèle.\par\smallskip\textbf{Démarche.}\enspace Le sens d’évolution compte : la température venait d’une valeur négative pendant un apport d’énergie.\par\smallskip{\small\textbf{Barème.}\enspace 0,5 pour fusion ; 0,5 pour coexistence solide-liquide.\par}\end{minipage}\par
\par\noindent\begin{minipage}{\linewidth}\pcSousTitre{Question 2}Le palier dure 6 − 2 = 4 min. La masse d’eau reste 200 g : la même eau est conservée sans fuite.\par\smallskip\textbf{Démarche.}\enspace Une durée est une différence de dates. La conservation de la masse concerne le système fermé, même quand son état change.\par\smallskip{\small\textbf{Barème.}\enspace 0,5 pour durée ; 0,5 pour masse justifiée.\par}\end{minipage}\par
\par\noindent\begin{minipage}{\linewidth}\pcSousTitre{Question 3}L’affirmation est fausse dans cette expérience. L’apport d’énergie se poursuit et permet le changement d’état ; pendant la fusion d’un corps pur à pression fixée, la température peut rester constante.\par\smallskip\textbf{Démarche.}\enspace Une température constante ne signifie pas l’absence de transfert d’énergie. Relier l’apport au changement d’état observé.\par\smallskip{\small\textbf{Barème.}\enspace 0,5 pour affirmation rejetée ; 0,5 pour énergie du changement d’état.\par}\end{minipage}\par
\pcFigureEx{figures/ann14-fusion-cor.pdf}{Modèle de correction · les légendes précisent les directions, les connexions ou les résultats.}
\pcPage{Corrigé du sujet 14}{Analyser un liquide de conditionnement}
\pcRepere{ann-14}{Revenir à l’énoncé}
\par\noindent\begin{minipage}{\linewidth}\pcSousTitre{Question 4}Il est basique : le pH 8,0 est supérieur à 7 vers 25 °C.\par\smallskip\textbf{Démarche.}\enspace Comparer le pH à 7 à la température précisée. Une valeur voisine de 7 n’est pas forcément exactement neutre.\par\smallskip{\small\textbf{Barème.}\enspace 0,25 pour basique ; 0,25 pour comparaison à 7.\par}\end{minipage}\par
\par\noindent\begin{minipage}{\linewidth}\pcSousTitre{Question 5}Le test met en évidence l’ion chlorure $\ce{Cl-}$. Avec 17 protons et une charge $-1$, il possède 18 électrons.\par\smallskip\textbf{Démarche.}\enspace Un ion négatif possède un excès d’électrons. La charge ne change pas le numéro atomique de l’élément.\par\smallskip{\small\textbf{Barème.}\enspace 0,5 pour identification ; 0,5 pour 18 électrons justifiés.\par}\end{minipage}\par
\par\noindent\begin{minipage}{\linewidth}\pcSousTitre{Question 6}Les résultats indiquent un liquide basique dans ces conditions et la présence d’ions chlorure dans la limite du test. Ils ne donnent pas tous les constituants ni leur quantité. Ils ne suffisent pas à garantir l’innocuité : il faut d’autres analyses et les informations adaptées sur le produit et son usage.\par\smallskip\textbf{Démarche.}\enspace Séparer trois conclusions : acidité/basicité, présence d’une espèce et évaluation d’un risque. Un test positif ne démontre pas la pureté ni l’absence d’autres substances.\par\smallskip{\small\textbf{Barème.}\enspace 0,5 pour les deux résultats exacts ; 0,5 pour composition non exhaustive ; 0,5 pour innocuité non démontrée.\par}\end{minipage}\par
\pcFigureEx{figures/ann14-ion.pdf}{Modèle de correction · les légendes précisent les directions, les connexions ou les résultats.}
\pcPage{Corrigé du sujet 14}{Limiter la vitesse sur un convoyeur}
\pcRepere{ann-14}{Revenir à l’énoncé}
\par\noindent\begin{minipage}{\linewidth}\pcSousTitre{Question 7}$d=v\Delta t=0{,}50\times8{,}0=4{,}0\unite m$.\par\smallskip\textbf{Démarche.}\enspace Multiplier la vitesse par la durée seulement pour la portion où la vitesse garde la valeur annoncée.\par\smallskip{\small\textbf{Barème.}\enspace 0,5 pour relation et produit ; 0,5 pour distance avec unité.\par}\end{minipage}\par
\par\noindent\begin{minipage}{\linewidth}\pcSousTitre{Question 8}À 0,50 m/s : $E_c=0{,}5\times2{,}0\times0{,}50^2=0{,}25\unite J$. À 1,0 m/s : $E_c=1{,}0\unite J$.\par\smallskip\textbf{Démarche.}\enspace Mettre toute la valeur de la vitesse au carré. Une masse en kilogrammes et une vitesse en m/s donnent une énergie en joules.\par\smallskip{\small\textbf{Barème.}\enspace 0,5 pour chaque calcul avec unité.\par}\end{minipage}\par
\par\noindent\begin{minipage}{\linewidth}\pcSousTitre{Question 9}Non : à masse constante, doubler la vitesse multiplie l’énergie cinétique par quatre. À vitesse plus faible, il y a moins d’énergie cinétique à transférer lors d’un arrêt ou d’un choc, toutes choses égales par ailleurs.\par\smallskip\textbf{Démarche.}\enspace Comparer les deux énergies calculées et relier le facteur quatre au carré de la vitesse.\par\smallskip{\small\textbf{Barème.}\enspace 0,5 pour facteur quatre ; 0,5 pour intérêt formulé en termes d’énergie.\par}\end{minipage}\par
\pcFigureEx{figures/ann14-energies.pdf}{Modèle de correction · les légendes précisent les directions, les connexions ou les résultats.}
\par\noindent\begin{minipage}{\linewidth}{\small\textbf{Bilan :} /3 + /3 + /3 + /1 = /10.\par\textbf{Pour reprendre :} \hyperref[chap-1]{chapitre 1 (p.~\pageref*{chap-1})} ; \hyperref[chap-2]{chapitre 2 (p.~\pageref*{chap-2})} ; \hyperref[chap-4]{chapitre 4 (p.~\pageref*{chap-4})} ; \hyperref[chap-8]{chapitre 8 (p.~\pageref*{chap-8})} ; \hyperref[chap-11]{chapitre 11 (p.~\pageref*{chap-11})}.\par}\end{minipage}\par
\pcPage{Corrigé du sujet 15}{Décrire le gaz de l’enveloppe}
\phantomsection\label{cor-ann-15}
\pcRepere{ann-15}{Revenir à l’énoncé}
\par\noindent\begin{minipage}{\linewidth}\pcSousTitre{Question 1}Il possède 2 protons, 4 − 2 = 2 neutrons et 2 électrons puisqu’il est neutre.\par\smallskip\textbf{Démarche.}\enspace Utiliser Z, puis A − Z, puis la neutralité électrique ; A ne compte pas les électrons.\par\smallskip{\small\textbf{Barème.}\enspace 0,25 pour protons ; 0,25 pour neutrons ; 0,5 pour électrons et neutralité.\par}\end{minipage}\par
\par\noindent\begin{minipage}{\linewidth}\pcSousTitre{Question 2}Non : les entités décrites sont des atomes d’hélium isolés. Il s’agit ici d’un gaz monoatomique.\par\smallskip\textbf{Démarche.}\enspace Une entité n’est pas nécessairement une molécule : tenir compte de la nature du gaz indiquée dans le document.\par\smallskip{\small\textbf{Barème.}\enspace 0,5 pour non ; 0,5 pour atomes isolés.\par}\end{minipage}\par
\par\noindent\begin{minipage}{\linewidth}\pcSousTitre{Question 3}Le nombre d’atomes ne diminue pas : aucun gaz n’a quitté le système. On représente les mêmes atomes, de même taille conventionnelle, dans un volume plus petit ; leurs espacements moyens diminuent.\par\smallskip\textbf{Démarche.}\enspace Distinguer la taille des entités, leur nombre et l’espace disponible. Une compression ne détruit pas des atomes.\par\smallskip{\small\textbf{Barème.}\enspace 0,5 pour conservation du nombre ; 0,5 pour réduction des espacements.\par}\end{minipage}\par
\pcFigureEx{figures/ann15-compression.pdf}{Modèle de correction · les légendes précisent les directions, les connexions ou les résultats.}
\pcPage{Corrigé du sujet 15}{Équilibrer la maquette}
\pcRepere{ann-15}{Revenir à l’énoncé}
\par\noindent\begin{minipage}{\linewidth}\pcSousTitre{Question 4}$P=mg=0{,}30\times10=3{,}0\unite N$.\par\smallskip\textbf{Démarche.}\enspace La masse est déjà en kilogrammes ; le produit par g donne un poids en newtons.\par\smallskip{\small\textbf{Barème.}\enspace 0,25 pour produit ; 0,25 pour résultat avec unité.\par}\end{minipage}\par
\par\noindent\begin{minipage}{\linewidth}\pcSousTitre{Question 5}Non : leurs valeurs sont différentes, 3,4 N vers le haut et 3,0 N vers le bas. La résultante du modèle vaut 0,4 N vers le haut.\par\smallskip\textbf{Démarche.}\enspace Deux forces opposées ne se compensent que si leurs valeurs sont égales.\par\smallskip{\small\textbf{Barème.}\enspace 0,5 pour comparaison des valeurs ; 0,5 pour résultante vers le haut.\par}\end{minipage}\par
\par\noindent\begin{minipage}{\linewidth}\pcSousTitre{Question 6}Le fil exerce une force de 0,4 N vers le bas. Avec le poids, la somme des forces vers le bas vaut 3,0 + 0,4 = 3,4 N, égale à la force de l’air vers le haut.\par\smallskip\textbf{Démarche.}\enspace Isoler le dirigeable comme système. Le fil ancré au sol tire la maquette vers le bas ; il ne la pousse pas vers le haut.\par\smallskip{\small\textbf{Barème.}\enspace 0,5 pour sens ; 0,5 pour 0,4 N justifié ; 0,5 pour schéma complet.\par}\end{minipage}\par
\pcFigureEx{figures/ann15-forces-cor.pdf}{Modèle de correction · les légendes précisent les directions, les connexions ou les résultats.}
\pcPage{Corrigé du sujet 15}{Alimenter les accessoires}
\pcRepere{ann-15}{Revenir à l’énoncé}
\par\noindent\begin{minipage}{\linewidth}\pcSousTitre{Question 7}Chaque branche reçoit 6,0 V. L’intensité totale est $I=0{,}20+0{,}10=0{,}30\unite A$.\par\smallskip\textbf{Démarche.}\enspace En dérivation, la tension est commune aux branches et le courant principal est la somme des courants de branche.\par\smallskip{\small\textbf{Barème.}\enspace 0,5 pour tension commune ; 0,5 pour intensité totale.\par}\end{minipage}\par
\par\noindent\begin{minipage}{\linewidth}\pcSousTitre{Question 8}$P=6{,}0\times0{,}30=1{,}8\unite W$ ; $\Delta t=E/P=1{,}8/1{,}8=1{,}0\unite h$.\par\smallskip\textbf{Démarche.}\enspace Une énergie en Wh divisée par une puissance en W donne une durée en heures. La capacité fournie est déjà l’énergie utilisable.\par\smallskip{\small\textbf{Barème.}\enspace 0,5 pour puissance ; 0,5 pour durée avec unité.\par}\end{minipage}\par
\par\noindent\begin{minipage}{\linewidth}\pcSousTitre{Question 9}Non, dans le modèle à générateur adapté et tension maintenue. La branche du moteur conserve une boucle fermée entre les bornes du générateur. Seule la branche d’éclairage est interrompue.\par\smallskip\textbf{Démarche.}\enspace Suivre un trajet complet incluant le générateur pour chaque branche. Une ouverture dans une branche ne coupe pas automatiquement les autres.\par\smallskip{\small\textbf{Barème.}\enspace 0,5 pour réponse ; 0,5 pour boucle du moteur conservée.\par}\end{minipage}\par
\pcFigureEx{figures/ann15-circuit-cor.pdf}{Modèle de correction · les légendes précisent les directions, les connexions ou les résultats.}
\par\noindent\begin{minipage}{\linewidth}{\small\textbf{Bilan :} /3 + /3 + /3 + /1 = /10.\par\textbf{Pour reprendre :} \hyperref[chap-2]{chapitre 2 (p.~\pageref*{chap-2})} ; \hyperref[chap-9]{chapitre 9 (p.~\pageref*{chap-9})} ; \hyperref[chap-13]{chapitre 13 (p.~\pageref*{chap-13})} ; \hyperref[chap-14]{chapitre 14 (p.~\pageref*{chap-14})} ; \hyperref[chap-15]{chapitre 15 (p.~\pageref*{chap-15})}.\par}\end{minipage}\par
\pcPage{Sources des sujets 11 à 15}{Des attendus étudiés, des documents et des problèmes nouveaux}
\phantomsection\label{sources-annales}
\begin{pcEncadre}{method}{Ce qui a été étudié}Les annales officielles servent à analyser la place des documents, les changements d’unité, le raisonnement progressif, les constructions et les réponses argumentées. Le volet PC du sujet 2025 consulté est noté sur 25, celui de 2026 sur 10. Pour ces entraînements 2027, les critères scientifiques et le point de langue suivent le cadre actuel et restent séparés.\end{pcEncadre}
\pcSousTitre{Sources primaires consultées le 24 octobre 2026}
\begin{questions}
\item Ministère de l’Éducation nationale, DNB 2025, sciences, série générale, sujet 25GENSCAN1 : analyse de la progression des questions, des documents et de la tâche de synthèse.
\item Ministère de l’Éducation nationale, DNB 2026, sciences, série générale, sujet 26GENSCAN1 : analyse de l’usage de mesures et de documents indépendants.
\item Éduscol, sujet de référence de physique-chimie pour le DNB 2027 : analyse des attendus de rédaction et des critères de réponse.
\item Bulletin officiel, note du 11 septembre 2026, NOR MENE2623228N, partie I.4 : durée et modalités de sciences applicables à compter de la session 2027.
\end{questions}
\pcSousTitre{Consulter les documents officiels}
{\small\url{https://eduscol.education.gouv.fr/5607/les-epreuves-du-dnb}\par
\url{https://www.education.gouv.fr/bo/2026/Hebdo4/MENE2623228N}\par
\url{https://www.education.gouv.fr/brevet-bac-et-cap-les-sujets-des-examens-2026-504911}\par}
\pcSousTitre{Des choix éditoriaux explicites}
Le monte-charge, le détecteur optique, la station météo, le colis réfrigéré et le dirigeable constituent cinq dossiers originaux. Les séries numériques et les schémas ont été construits pour ces exercices. Aucun sujet n’annonce les thèmes qui seront retenus à l’examen. La concordance générale du programme figure dans les annexes du manuel.
\pcAnnexe{Méthodes et références utiles}{Des fiches à consulter puis à utiliser sans aide}\phantomsection\label{methodes}\pdfbookmark[0]{Méthodes et outils}{bm-methodes}
\begin{pcEncadre}{method}{Calculer en cinq étapes}1. Nommer la grandeur cherchée. 2. Écrire une relation adaptée. 3. Convertir les données. 4. Remplacer les valeurs et calculer. 5. Donner l’unité, un arrondi utile et une conclusion dans le contexte.\end{pcEncadre}
\pcSousTitre{Exemple : une énergie de fonctionnement}
Un appareil de 1,2 kW fonctionne 15 min à puissance constante. En kWh, écrire $\Delta t=15/60=0{,}25\unite h$, puis $E=1{,}2\times0{,}25=0{,}30\unite{kWh}$. En joules, utiliser 1 200 W et 900 s : $E=1\,080\,000\unite J$. Les deux résultats sont cohérents.
\pcSousTitre{Isoler une grandeur}
De $P=mg$ : $m=P/g$ et $g=P/m$, si les dénominateurs sont non nuls. De $E=P\Delta t$ : $P=E/\Delta t$ et $\Delta t=E/P$. Vérifier le sens : à même énergie, une puissance deux fois plus faible demande deux fois plus de temps.
\pcSousTitre{Comparer avant de calculer}
Si $E_c=\tfrac12mv^2$, doubler m double Ec ; doubler v quadruple Ec. Si la gravitation est fournie sous $F=Gm_Am_B/d^2$, doubler d divise F par quatre. Énoncer ce qui reste fixé.
\pcSousTitre{Contrôler un ordre de grandeur}
Une vitesse de vélo de plusieurs milliers de m/s ou une énergie en N signale une erreur probable de conversion ou de grandeur. Le contrôle de plausibilité complète le calcul ; il ne remplace pas les relations.
\pcPage{Les relations de fin de cycle}{Grandeurs, unités et conditions d’emploi}
\phantomsection\label{formules}
\begingroup\renewcommand{\arraystretch}{1.5}\fontsize{10.5}{13}\selectfont
\begin{tabularx}{\linewidth}{@{}p{37mm}p{37mm}Y@{}}\toprule
Relation & Unités cohérentes & Conditions et précautions\\\midrule
$\rho=m/V$ & kg ; m³ ; kg/m³ & Matériau et température précisés\\
$v=d/\Delta t$ & m ; s ; m/s & Moyenne sur le trajet ; uniforme si précisé\\
$P=mg$ & kg ; N/kg ; N & g du lieu ; P est ici le poids\\
$F=Gm_Am_B/d^2$ & kg ; m ; N & Loi fournie ; d entre centres ; G en SI\\
$E_c=\tfrac12mv^2$ & kg ; m/s ; J & Référentiel et système indiqués\\
$U=RI$ & V ; A ; $\Omega$ & Conducteur ohmique, conditions fixées\\
$P=UI$ & V ; A ; W & Courant continu dans le cours\\
$E=P\Delta t$ & W ; s ; J & Puissance constante ou moyenne fournie\\
$E=P\Delta t$ & kW ; h ; kWh & Ne pas mélanger W et kW\\
$d=v\Delta t$ & m ; s ; m/s & Distance totale parcourue par le signal\\
$f=N/\Delta t$ & répétitions ; s ; Hz & Signal périodique ; N motifs complets\\\bottomrule
\end{tabularx}\endgroup
\begin{pcEncadre}{attention}{Le contexte détermine les symboles}
P peut désigner un poids en N ou une puissance en W. Lire la légende et l’unité, puis écrire le nom de la grandeur au début de la démarche.
\end{pcEncadre}
\pcSousTitre{Lois des circuits}
En série dans une branche : même I. À un nœud : somme entrante = somme sortante. Dans une boucle source + deux récepteurs avec conventions adaptées : $U_G=U_1+U_2$. En dérivation entre les mêmes nœuds : même U.
\pcPage{Conversions et puissances de dix}{Une fiche pour éviter les erreurs d’unité}
\phantomsection\label{conversions}
\begin{questions}
\item $1\unite{kg}=1\,000\unite g$ ; $1\unite g=10^{-3}\unite{kg}$.
\item $1\unite L=1\unite{dm^3}$ ; $1\unite{mL}=1\unite{cm^3}$ ; $1\unite{m^3}=1\,000\unite L$.
\item $1\unite{g/mL}=1\,000\unite{kg/m^3}$. Un cube de 1 m de côté a un volume de $10^6\unite{cm^3}$.
\item $1\unite h=3\,600\unite s$ ; $1\unite{min}=60\unite s$. 15 min = 0,25 h, et non 0,15 h.
\item De km/h vers m/s : diviser par 3,6. De m/s vers km/h : multiplier par 3,6.
\item $1\unite A=1\,000\unite{mA}$ ; $1\unite{k\Omega}=1\,000\unite\Omega$ ; $1\unite{kW}=1\,000\unite W$.
\item $1\unite{Wh}=3\,600\unite J$ ; $1\unite{kWh}=3{,}6\times10^6\unite J$. Le kWh est une unité d’énergie.
\item $1\unite{kHz}=1\,000\unite{Hz}$ ; $1\unite{al}\approx9{,}46\times10^{15}\unite m$.
\end{questions}
\begin{pcEncadre}{method}{Le carré concerne toute la valeur}$(2\times10^3)^2=4\times10^6$ et non $2\times10^6$. Sur la calculatrice, saisir les parenthèses autour d’une distance écrite avec une puissance de dix avant de la mettre au carré.\end{pcEncadre}
\pcPage{Lire un graphique}{Exploiter un document et respecter la précision des données}
\phantomsection\label{graphiques}
\begin{questions}
\item Lire le titre, les noms des axes, les unités, l’origine et les graduations. Vérifier si l’axe des intensités est en A ou en mA.
\item Repérer la série et la légende ; distinguer mesure, modèle tracé et valeur de référence.
\item Pour une lecture, projeter le point vers les axes. Donner une valeur arrondie compatible avec la graduation.
\item Une droite ne prouve une proportionnalité que si elle est compatible avec le passage par l’origine et les conditions étudiées.
\item Pour un tableau, citer au moins une valeur pertinente et son unité avant de comparer.
\item Identifier l’auteur d’un document, sa date, son périmètre et les éventuelles données fictives. Une image seule n’est pas une mesure.
\end{questions}
\begin{pcEncadre}{method}{Construire une réponse ouverte}Dire ce que l’on cherche ; choisir les données utiles ; expliquer la relation ; calculer si nécessaire ; comparer ; conclure avec une limite. Si deux options ont des avantages différents, préciser le critère plutôt que proclamer une supériorité absolue.\end{pcEncadre}
\pcSousTitre{Exemple de réponse scientifique courte}
« Dans le modèle fourni, A émet 1 kg de CO2 et B en émet 4 kg. A a donc des émissions quatre fois plus faibles pour ce trajet. Cette comparaison porte seulement sur le périmètre et l’indicateur fournis ; elle ne décrit pas tous les impacts. »\pcPage{Air, solutions et climat}{Révisions transversales du cycle 4}
\phantomsection\label{air}
\pcSousTitre{Composition et pression}
L’air sec contient environ 78 \% de diazote, 21 \% de dioxygène et 1 \% d’autres gaz en volume. La vapeur d’eau est variable. Un gaz possède une masse et se comprime : à quantité de gaz conservée et après retour à la même température, un volume plus petit correspond à une pression plus élevée. La pression se mesure en pascals, Pa ; $1\unite{hPa}=100\unite{Pa}$ et $1\unite{bar}=10^5\unite{Pa}$. La pression atmosphérique usuelle est voisine de $10^5\unite{Pa}$.
\pcSousTitre{Dissolution et traitement des eaux}
Des gaz comme CO2 peuvent se dissoudre dans l’eau. Un gaz qui se dégage pouvait être déjà dissous : le dégazage n’est pas toujours une réaction. La miscibilité concerne deux liquides ; la solubilité caractérise une limite de dissolution dans des conditions données. Décantation, filtration et distillation séparent selon des propriétés différentes. Une eau clarifiée n’est pas automatiquement potable.
\pcSousTitre{Gaz à effet de serre}
CO2, vapeur d’eau, méthane et protoxyde d’azote font partie des gaz à effet de serre. Ils absorbent une partie du rayonnement infrarouge terrestre et émettent aussi du rayonnement. Un déséquilibre durable entre énergie reçue et énergie renvoyée peut modifier l’énergie du système Terre-atmosphère et sa température. L’effet de serre naturel est distinct de son renforcement lié aux activités humaines.
\pcSousTitre{Des conclusions qui respectent le périmètre}
Une absence de CO2 au point d’usage ne prouve pas une absence d’émissions sur tout le cycle de vie. L’acidification est une baisse de pH : un océan peut devenir moins basique tout en gardant un pH supérieur à 7. Comparer les émissions exige une même quantité de service et un périmètre annoncé.
\begin{pcEncadre}{method}{Trois questions flash, corrigées}
Dans 100 L d’air sec : environ 21 L de dioxygène. Une solution pH 8,0 reste basique vers 25 °C. Un écart énergétique reçu − renvoyé de 10 J correspond, dans un bilan sans autre échange, à une hausse d’énergie du système de 10 J.
\end{pcEncadre}
\pcPage{Tests et formules usuelles}{Une banque à exploiter avec le protocole donné}
\phantomsection\label{tests}
\begingroup\renewcommand{\arraystretch}{1.45}\begin{tabularx}{\linewidth}{@{}p{32mm}p{47mm}Y@{}}\toprule
Espèce recherchée & Test fourni & Résultat positif usuel\\\midrule
CO2 & Eau de chaux & Trouble\\
O2 & Bûchette incandescente & Ravivement, test encadré\\
H2 & Flamme, test sécurisé & Petite détonation, test encadré\\
Eau & Sulfate de cuivre anhydre & Blanc vers bleu\\
Cl− & Nitrate d’argent & Précipité blanc\\
Cu2+ & Soude & Précipité bleu\\
Fe2+ & Soude & Précipité vert\\
Fe3+ & Soude & Précipité brun\\\bottomrule
\end{tabularx}\endgroup
\pcSousTitre{Les formules à reconnaître}
Dihydrogène H2 ; dioxygène O2 ; diazote N2 ; eau H2O ; dioxyde de carbone CO2 ; méthane CH4 ; protoxyde d’azote N2O. Ces écritures décrivent la composition ; les noms ne se remplacent pas par « gaz » sans précision.
\begin{pcEncadre}{attention}{Une banque n’est pas une consigne de manipulation libre}
Utiliser les quantités, protections et traitements des déchets indiqués par le professeur. Un test positif s’interprète dans ses conditions ; un résultat négatif dépend aussi de la sensibilité et de l’état du réactif. Faire les tests sur des prélèvements séparés.
\end{pcEncadre}
\pcPage{Repères des éléments}{Z identifie le nombre de protons}
\phantomsection\label{elements}
Sélection de repérage ; ce tableau ne remplace pas une classification périodique complète. A n’est pas donné ici : on ne déduit pas un nombre précis de neutrons du seul Z.\par\vspace{3mm}
\begingroup\fontsize{10.5}{13}\selectfont\renewcommand{\arraystretch}{1.12}\begin{tabularx}{.85\linewidth}{@{}p{20mm}Yp{25mm}@{}}\toprule Z & Élément & Symbole\\\midrule
1 & Hydrogène & H\\
2 & Hélium & He\\
3 & Lithium & Li\\
4 & Béryllium & Be\\
5 & Bore & B\\
6 & Carbone & C\\
7 & Azote & N\\
8 & Oxygène & O\\
9 & Fluor & F\\
10 & Néon & Ne\\
11 & Sodium & Na\\
12 & Magnésium & Mg\\
13 & Aluminium & Al\\
14 & Silicium & Si\\
15 & Phosphore & P\\
16 & Soufre & S\\
17 & Chlore & Cl\\
18 & Argon & Ar\\
19 & Potassium & K\\
20 & Calcium & Ca\\
26 & Fer & Fe\\
29 & Cuivre & Cu\\
30 & Zinc & Zn\\
\bottomrule\end{tabularx}\endgroup
\pcPage{Diagnostic : corrigé}{Un repère pour reprendre le bon chapitre}
\phantomsection\label{cor-diagnostic}
\begin{questions}
\item 250 mL et 0,350 kg. Reprendre la fiche conversions.
\item Non : une solution homogène peut contenir plusieurs espèces ; chapitre 1.
\item Quatre H et deux O ; chapitre 2.
\item 11 électrons, par neutralité ; chapitre 2.
\item Acide, car pH < 7 ; chapitre 4.
\item $v=120/30=4\unite{m/s}$ ; chapitre 8.
\item Poids en N ; masse en kg ; chapitre 9.
\item Ampèremètre en série ; voltmètre en dérivation ; chapitre 13.
\item $0{,}020\unite A$ et $900\unite s$ ; conversions.
\item La distance totale comprend l’aller et le retour ; chapitre 17.
\end{questions}
\begin{pcEncadre}{method}{Décider de la suite}Si une conversion est erronée, refaire deux exemples avec leurs unités. Si la notion est confuse, relire le cours et expliquer le schéma. Si la réponse est juste mais difficile à suivre, retravailler les phrases et l’organisation.\end{pcEncadre}
\pcPage{Lexique scientifique}{Des définitions à replacer dans une situation}
\phantomsection\label{lexique}
\begingroup\fontsize{10.5}{13.5}\selectfont\begin{multicols}{2}
\par\addvspace{2mm}\noindent\begin{minipage}{\linewidth}\textbf{Atome}\enspace Entité neutre comportant un noyau et des électrons.\par\end{minipage}\par
\par\addvspace{2mm}\noindent\begin{minipage}{\linewidth}\textbf{Ion}\enspace Entité chargée ayant perdu ou gagné des électrons.\par\end{minipage}\par
\par\addvspace{2mm}\noindent\begin{minipage}{\linewidth}\textbf{Nucléon}\enspace Proton ou neutron du noyau.\par\end{minipage}\par
\par\addvspace{2mm}\noindent\begin{minipage}{\linewidth}\textbf{Numéro atomique}\enspace Z : nombre de protons ; identifie l’élément.\par\end{minipage}\par
\par\addvspace{2mm}\noindent\begin{minipage}{\linewidth}\textbf{Molécule}\enspace Assemblage neutre d’atomes liés.\par\end{minipage}\par
\par\addvspace{2mm}\noindent\begin{minipage}{\linewidth}\textbf{Espèce chimique}\enspace Ensemble d’entités chimiques identiques.\par\end{minipage}\par
\par\addvspace{2mm}\noindent\begin{minipage}{\linewidth}\textbf{Réactif}\enspace Espèce consommée lors d’une réaction.\par\end{minipage}\par
\par\addvspace{2mm}\noindent\begin{minipage}{\linewidth}\textbf{Produit}\enspace Espèce formée lors d’une réaction.\par\end{minipage}\par
\par\addvspace{2mm}\noindent\begin{minipage}{\linewidth}\textbf{Corrosion}\enspace Transformation d’un métal avec son environnement.\par\end{minipage}\par
\par\addvspace{2mm}\noindent\begin{minipage}{\linewidth}\textbf{Dissolution}\enspace Dispersion d’un soluté dans un solvant.\par\end{minipage}\par
\par\addvspace{2mm}\noindent\begin{minipage}{\linewidth}\textbf{Dilution}\enspace Ajout de solvant sans ajout de soluté.\par\end{minipage}\par
\par\addvspace{2mm}\noindent\begin{minipage}{\linewidth}\textbf{Masse volumique}\enspace Masse divisée par le volume, dans des conditions précisées.\par\end{minipage}\par
\par\addvspace{2mm}\noindent\begin{minipage}{\linewidth}\textbf{Force}\enspace Modèle d’une action : point, direction, sens et valeur.\par\end{minipage}\par
\par\addvspace{2mm}\noindent\begin{minipage}{\linewidth}\textbf{Poids}\enspace Force de pesanteur exercée sur un objet.\par\end{minipage}\par
\par\addvspace{2mm}\noindent\begin{minipage}{\linewidth}\textbf{Référentiel}\enspace Référence choisie pour décrire les positions et leur évolution.\par\end{minipage}\par
\par\addvspace{2mm}\noindent\begin{minipage}{\linewidth}\textbf{Trajectoire}\enspace Ensemble des positions successives d’un point.\par\end{minipage}\par
\par\addvspace{2mm}\noindent\begin{minipage}{\linewidth}\textbf{Énergie cinétique}\enspace Énergie associée au mouvement dans un référentiel.\par\end{minipage}\par
\par\addvspace{2mm}\noindent\begin{minipage}{\linewidth}\textbf{Transfert}\enspace Passage d’énergie entre systèmes.\par\end{minipage}\par
\par\addvspace{2mm}\noindent\begin{minipage}{\linewidth}\textbf{Conversion}\enspace Changement de forme d’énergie.\par\end{minipage}\par
\par\addvspace{2mm}\noindent\begin{minipage}{\linewidth}\textbf{Puissance}\enspace Énergie transférée par unité de durée.\par\end{minipage}\par
\par\addvspace{2mm}\noindent\begin{minipage}{\linewidth}\textbf{Tension}\enspace Grandeur électrique entre deux points, en V.\par\end{minipage}\par
\par\addvspace{2mm}\noindent\begin{minipage}{\linewidth}\textbf{Intensité}\enspace Grandeur caractérisant le courant, en A.\par\end{minipage}\par
\par\addvspace{2mm}\noindent\begin{minipage}{\linewidth}\textbf{Résistance}\enspace Grandeur d’un conducteur ohmique, en Ω.\par\end{minipage}\par
\par\addvspace{2mm}\noindent\begin{minipage}{\linewidth}\textbf{Fréquence}\enspace Nombre de répétitions périodiques par seconde.\par\end{minipage}\par
\par\addvspace{2mm}\noindent\begin{minipage}{\linewidth}\textbf{Signal}\enspace Phénomène détectable pouvant transporter de l’information.\par\end{minipage}\par
\par\addvspace{2mm}\noindent\begin{minipage}{\linewidth}\textbf{Année-lumière}\enspace Distance parcourue par la lumière dans le vide en un an.\par\end{minipage}\par
\par\addvspace{2mm}\noindent\begin{minipage}{\linewidth}\textbf{Témoin}\enspace Comparaison permettant d’examiner une cause possible.\par\end{minipage}\par
\par\addvspace{2mm}\noindent\begin{minipage}{\linewidth}\textbf{Modèle}\enspace Représentation simplifiée dont le domaine d’emploi est précisé.\par\end{minipage}\par
\end{multicols}\endgroup
\pcPage{Programme et concordance}{Repérer les apprentissages du cycle 4}
\phantomsection\label{programme}
\pdfbookmark[0]{Programme et références}{bm-programme}
Les lignes relient les notions du programme en vigueur aux cours, exercices et sujets blancs. Le découpage annuel reste éditorial. La loi de gravitation est utilisée quand elle est fournie ; la formule de l’énergie potentielle n’est pas un prérequis exigé par ce manuel.\par\vspace{2mm}
\par\noindent\begin{minipage}{\linewidth}\textbf{M01} · États, changements d’état et interprétation microscopique. \hyperref[chap-1]{Chapitre 1}, exercices 1, 3, 10, 14. \emph{Cours ; sujet 7.}\par\end{minipage}\par\vspace{2mm}
\par\noindent\begin{minipage}{\linewidth}\textbf{M02} · Corps purs, mélanges, espèces et séparation. \hyperref[chap-1]{Chapitre 1}, exercices 2, 9, 11, 15. \emph{Fiche air ; sujet 2.}\par\end{minipage}\par\vspace{2mm}
\par\noindent\begin{minipage}{\linewidth}\textbf{M03} · Masse volumique, mesure, unités et température. \hyperref[chap-1]{Chapitre 1}, exercices 5, 6, 7, 8, 12, 13. \emph{Sujets 2 et 7.}\par\end{minipage}\par\vspace{2mm}
\par\noindent\begin{minipage}{\linewidth}\textbf{M04} · Solubilité, miscibilité et dissolution de gaz. \hyperref[chap-1]{Chapitre 1}, exercices 9, 11, 15. \emph{Fiche air, solutions et climat.}\par\end{minipage}\par\vspace{2mm}
\par\noindent\begin{minipage}{\linewidth}\textbf{M05} · Atomes, molécules, ions, formules et tableau périodique. \hyperref[chap-2]{Chapitre 2}, exercices 4, 6, 7, 8, 9, 14. \emph{Sujets 2, 5 et 7.}\par\end{minipage}\par\vspace{2mm}
\par\noindent\begin{minipage}{\linewidth}\textbf{M06} · Protons, neutrons, électrons et structure du noyau. \hyperref[chap-2]{Chapitre 2}, exercices 1, 2, 3, 8, 13, 15. \emph{Fiche des éléments ; sujet 7.}\par\end{minipage}\par\vspace{2mm}
\par\noindent\begin{minipage}{\linewidth}\textbf{M07} · Tests caractéristiques et transformations physiques/chimiques. \hyperref[chap-3]{Chapitre 3}, exercices 1, 3, 11, 12, 14. \emph{Sujets 5 et 10.}\par\end{minipage}\par\vspace{2mm}
\par\noindent\begin{minipage}{\linewidth}\textbf{M08} · Équation fournie, redistribution des atomes et conservation. \hyperref[chap-3]{Chapitre 3}, exercices 4, 5, 7, 8, 9, 10, 13. \emph{Sujets 1 et 10.}\par\end{minipage}\par\vspace{2mm}
\par\noindent\begin{minipage}{\linewidth}\textbf{M09} · Acidité, basicité, pH et ions H+ ou OH−. \hyperref[chap-4]{Chapitre 4}, exercices 1, 3, 6, 7, 10, 14. \emph{Sujets 2 et 5.}\par\end{minipage}\par\vspace{2mm}
\par\noindent\begin{minipage}{\linewidth}\textbf{M10} · Réactions acide-base et acide-métal. \hyperref[chap-5]{Chapitre 5}, exercices 2, 3, 6, 7, 12, 14. \emph{Réaction acide-base au chapitre 4.}\par\end{minipage}\par\vspace{2mm}
\par\noindent\begin{minipage}{\linewidth}\textbf{M11} · Corrosion, environnement et protections. \hyperref[chap-5]{Chapitre 5}, exercices 1, 4, 9, 10, 13, 15. \emph{Sujet 5.}\par\end{minipage}\par\vspace{2mm}
\par\noindent\begin{minipage}{\linewidth}\textbf{M12} · Air, combustions, gaz à effet de serre et impacts. \hyperref[chap-6]{Chapitre 6}, exercices 1, 6, 8, 11, 13, 14. \emph{Fiches air/climat ; sujet 1.}\par\end{minipage}\par\vspace{2mm}
\par\noindent\begin{minipage}{\linewidth}\textbf{M13} · Univers, Système solaire, évolution, éléments et ressources. \hyperref[chap-7]{Chapitre 7}, exercices 1, 2, 7, 8, 12, 15. \emph{Année-lumière au chapitre 16 ; sujet 9.}\par\end{minipage}\par\vspace{2mm}
\par\noindent\begin{minipage}{\linewidth}\textbf{V01} · Référentiel, relativité du mouvement et trajectoires. \hyperref[chap-8]{Chapitre 8}, exercices 1, 3, 8, 9, 13, 14. \emph{Sujet 1.}\par\end{minipage}\par\vspace{2mm}
\par\noindent\begin{minipage}{\linewidth}\textbf{V02} · Vitesse : valeur, direction, sens ; distance et durée. \hyperref[chap-8]{Chapitre 8}, exercices 4, 5, 6, 7, 10, 11, 12. \emph{Sujets 1 et 6.}\par\end{minipage}\par\vspace{2mm}
\par\noindent\begin{minipage}{\linewidth}\textbf{V03} · Interactions, forces et quatre caractéristiques. \hyperref[chap-9]{Chapitre 9}, exercices 1, 4, 9, 10, 13, 14. \emph{Sujet 4.}\par\end{minipage}\par\vspace{2mm}
\par\noindent\begin{minipage}{\linewidth}\textbf{V04} · Poids, masse et P = mg. \hyperref[chap-9]{Chapitre 9}, exercices 5, 6, 7, 8, 11, 15. \emph{Sujet 4.}\par\end{minipage}\par\vspace{2mm}
\par\noindent\begin{minipage}{\linewidth}\textbf{V05} · Exploitation de la loi fournie de gravitation. \hyperref[chap-10]{Chapitre 10}, exercices 5, 6, 8, 9, 10, 11, 14. \emph{Sujets 4 et 9.}\par\end{minipage}\par\vspace{2mm}
\par\noindent\begin{minipage}{\linewidth}\textbf{E01} · Énergie cinétique et énergie potentielle de position. \hyperref[chap-11]{Chapitre 11}, exercices 4, 6, 7, 8, 9, 15. \emph{Sujet 1.}\par\end{minipage}\par\vspace{2mm}
\par\noindent\begin{minipage}{\linewidth}\textbf{E02} · Sources, formes, transferts, conversions et conservation. \hyperref[chap-12]{Chapitre 12}, exercices 1, 2, 6, 7, 8, 15. \emph{Sujets 3, 7 et 8.}\par\end{minipage}\par\vspace{2mm}
\par\noindent\begin{minipage}{\linewidth}\textbf{E03} · Rayonnement, absorption et système Terre-atmosphère. \hyperref[chap-12]{Chapitre 12}, exercices 5, 9, 10, 14, 15. \emph{Fiche climat ; sujet 8.}\par\end{minipage}\par\vspace{2mm}
\par\noindent\begin{minipage}{\linewidth}\textbf{E04} · Circuits série/dérivation et cahier des charges. \hyperref[chap-13]{Chapitre 13}, exercices 3, 4, 8, 10, 13, 15. \emph{Sujet 3.}\par\end{minipage}\par\vspace{2mm}
\par\noindent\begin{minipage}{\linewidth}\textbf{E05} · Mesures, unicité et additivité des intensités/tensions. \hyperref[chap-13]{Chapitre 13}, exercices 1, 6, 7, 8, 11, 12. \emph{Sujet 3.}\par\end{minipage}\par\vspace{2mm}
\par\noindent\begin{minipage}{\linewidth}\textbf{E06} · Loi d’Ohm et caractéristique. \hyperref[chap-14]{Chapitre 14}, exercices 3, 4, 5, 8, 9, 13. \emph{Sujet 10.}\par\end{minipage}\par\vspace{2mm}
\par\noindent\begin{minipage}{\linewidth}\textbf{E07} · Puissance électrique P = UI. \hyperref[chap-14]{Chapitre 14}, exercices 6, 7, 9, 14, 15. \emph{Sujets 3, 6 et 10.}\par\end{minipage}\par\vspace{2mm}
\par\noindent\begin{minipage}{\linewidth}\textbf{E08} · Énergie, puissance, durée, consommation et stockage. \hyperref[chap-15]{Chapitre 15}, exercices 2, 3, 4, 6, 7, 11, 15. \emph{Sujets 3, 6 et 8.}\par\end{minipage}\par\vspace{2mm}
\par\noindent\begin{minipage}{\linewidth}\textbf{S01} · Sources, diffusion, rayon lumineux et propagation. \hyperref[chap-16]{Chapitre 16}, exercices 1, 3, 5, 8, 10, 11. \emph{Sujet 9.}\par\end{minipage}\par\vspace{2mm}
\par\noindent\begin{minipage}{\linewidth}\textbf{S02} · Vitesse de lumière, distances et année-lumière. \hyperref[chap-16]{Chapitre 16}, exercices 4, 6, 7, 12, 13, 15. \emph{Sujets 4 et 9.}\par\end{minipage}\par\vspace{2mm}
\par\noindent\begin{minipage}{\linewidth}\textbf{S03} · Propagation sonore, fréquence et plages audibles. \hyperref[chap-17]{Chapitre 17}, exercices 1, 2, 4, 6, 7, 13. \emph{Sujet 6.}\par\end{minipage}\par\vspace{2mm}
\par\noindent\begin{minipage}{\linewidth}\textbf{S04} · Écho, sonar, distance et prévention auditive. \hyperref[chap-17]{Chapitre 17}, exercices 3, 8, 9, 11, 14, 15. \emph{Sujet 6.}\par\end{minipage}\par\vspace{2mm}
\par\noindent\begin{minipage}{\linewidth}\textbf{S05} · Signaux, information, codage et mesures. \hyperref[chap-18]{Chapitre 18}, exercices 1, 4, 6, 7, 10, 12, 15. \emph{Sujets 4 et 6.}\par\end{minipage}\par\vspace{2mm}
\pcPage{Compétences et numérique}{Pratiquer les démarches, les langages et les outils}
\begin{questions}
\item Démarche expérimentale : hypothèse, témoin, répétition et conclusion, chapitres 1, 3, 4 et 5.
\item Concevoir un dispositif : protocole de mesure et montage répondant à une demande, chapitres 9, 13 et 14.
\item Mobiliser le numérique : nuages de points au tableur, chapitres 1, 10, 11 et 14 ; vidéo au chapitre 8 ; signaux enregistrés aux chapitres 17 et 18.
\item Pratiquer les langages : équations, forces, circuits et graphiques ; réponses rédigées de tous les sujets blancs.
\item Adopter une attitude responsable : produits, déchets et ressources ; sécurité électrique ; protection visuelle et auditive.
\item Se situer dans le temps et l’espace : origine des éléments, évolution de l’Univers et ordres de grandeur, chapitres 2, 7 et 16.
\end{questions}
\begin{pcEncadre}{method}{Atelier tableur reproductible}Saisir v = 0, 2, 4, 6, 8 m/s ; fixer m = 2 kg. Calculer Ec par une formule recopiable. Conserver les données d’origine. Produire un nuage de points Ec(v), puis Ec(v²), avec noms des axes et unités. Expliquer la différence des deux graphes.\end{pcEncadre}
\pcPage{Références et choix éditoriaux}{Sources primaires vérifiées le 24 octobre 2026}
\phantomsection\label{references}
\par\noindent\begin{minipage}{\linewidth}\pcSousTitre{Programme en application}
Éduscol, programmes du cycle 4 pour 2026-2027 ; physique-chimie : BO du 30 juillet 2020.\par{\fontsize{8.5}{11}\selectfont\raggedright\url{https://eduscol.education.gouv.fr/4362/enseigner-au-cycle-4}\par}\end{minipage}\par
\par\noindent\begin{minipage}{\linewidth}\pcSousTitre{Texte de référence}
Programme de physique-chimie du cycle 4, pages 96 à 106 ; découpage annuel éditorial.\par{\fontsize{8.5}{11}\selectfont\raggedright\url{https://eduscol.education.fr/document/621/download}\par}\end{minipage}\par
\par\noindent\begin{minipage}{\linewidth}\pcSousTitre{Épreuves du DNB 2027}
Note de service du 11 septembre 2026, NOR MENE2623228N, partie I.4 : sciences.\par{\fontsize{8.5}{11}\selectfont\raggedright\url{https://www.education.gouv.fr/bo/2026/Hebdo4/MENE2623228N}\par}\end{minipage}\par
\par\noindent\begin{minipage}{\linewidth}\pcSousTitre{Présentation des épreuves}
Éduscol : durée, notation et compétences ; page actualisée en octobre 2026.\par{\fontsize{8.5}{11}\selectfont\raggedright\url{https://eduscol.education.gouv.fr/5607/les-epreuves-du-dnb}\par}\end{minipage}\par
\par\noindent\begin{minipage}{\linewidth}\pcSousTitre{Sujet de référence}
Éduscol, physique-chimie, série générale, n° 01 : repérage des attendus et de la rédaction.\par{\fontsize{8.5}{11}\selectfont\raggedright\url{https://eduscol.education.gouv.fr/sites/default/files/document/dnb-sujet-de-reference-physique-chimie-serie-generale-n01-129673.pdf}\par}\end{minipage}\par
\pcPage{Références complémentaires}{Poursuivre la lecture des sources primaires}
\par\noindent\begin{minipage}{\linewidth}\pcSousTitre{Annale officielle 2026}
Ministère, sciences, série générale : lecture critique du format et des démarches attendues. Les énoncés ne sont pas reproduits.\par{\fontsize{8.5}{11}\selectfont\raggedright\url{https://www.education.gouv.fr/sites/default/files/document/diplome-national-du-brevet-2026-sciences-serie-generale-517070.pdf}\par}\end{minipage}\par
\par\noindent\begin{minipage}{\linewidth}\pcSousTitre{Atome et noyau}
CEA : protons, neutrons, électrons et numéro atomique.\par{\fontsize{8.5}{11}\selectfont\raggedright\url{https://www.cea.fr/comprendre/enseignants/pages/ressources-pedagogiques/essentiel-sur/physique-chimie/essentiel-sur-noyaux-des-atomes.aspx}\par}\end{minipage}\par
\par\noindent\begin{minipage}{\linewidth}\pcSousTitre{Unités}
BIPM, Brochure sur le SI : unités, symboles et préfixes.\par{\fontsize{8.5}{11}\selectfont\raggedright\url{https://www.bipm.org/fr/publications/si-brochure}\par}\end{minipage}\par
\par\noindent\begin{minipage}{\linewidth}\pcSousTitre{Chronologie cosmique}
NASA Science : âges approximatifs de l’Univers et du Système solaire.\par{\fontsize{8.5}{11}\selectfont\raggedright\url{https://science.nasa.gov/exoplanets/what-is-the-universe/}\par}\end{minipage}\par
\par\noindent\begin{minipage}{\linewidth}\pcSousTitre{Ressources pédagogiques}
Éduscol : accompagnement du programme de physique-chimie au cycle 4.\par{\fontsize{8.5}{11}\selectfont\raggedright\url{https://eduscol.education.gouv.fr/5739/ressources-d-accompagnement-du-programme-de-physique-chimie-au-cycle-4}\par}\end{minipage}\par
\pcPage{Références complémentaires}{Poursuivre la lecture des sources primaires}
\par\noindent\begin{minipage}{\linewidth}\pcSousTitre{Prévention électrique}
INRS : risques et prévention en électricité.\par{\fontsize{8.5}{11}\selectfont\raggedright\url{https://www.inrs.fr/risques/electriques/risques-electricite.html}\par}\end{minipage}\par
\par\noindent\begin{minipage}{\linewidth}\pcSousTitre{Éléments et symboles}
IUPAC : tableau périodique des éléments.\par{\fontsize{8.5}{11}\selectfont\raggedright\url{https://iupac.org/what-we-do/periodic-table-of-elements/}\par}\end{minipage}\par
\par\noindent\begin{minipage}{\linewidth}\pcSousTitre{Système solaire}
Agence spatiale canadienne : formation des systèmes solaires.\par{\fontsize{8.5}{11}\selectfont\raggedright\url{https://www.asc-csa.gc.ca/fra/blogue/2021/08/06/comment-se-forment-les-systemes-solaires.asp}\par}\end{minipage}\par
\pcSousTitre{Création et réutilisation}
Les sujets blancs, documents d’étude, barèmes et schémas sont des créations éditoriales OptiTAB. Les sujets ministériels ont été consultés pour analyser les attendus, sans reproduction de leurs exercices ni de leurs images. Les activités de révision sont harmonisées avec les autres manuels de la collection. Les nombres d’étude, prix et scénarios sont fictifs sauf indication explicite.
